\documentclass[10pt,a4paper]{article}
\usepackage[a4paper,margin=0.82in]{geometry}
\usepackage[T1]{fontenc}
\usepackage{lmodern,microtype}
\usepackage{amsmath,amssymb,amsthm,mathtools}
\usepackage{booktabs,array,tabularx,longtable,multirow,float}
\usepackage{enumitem,xcolor,url,listings}
\usepackage[hidelinks]{hyperref}
\usepackage[nameinlink,noabbrev,capitalise]{cleveref}
\usepackage{fancyhdr}

\definecolor{deepblue}{RGB}{28,69,125}
\hypersetup{colorlinks=true,linkcolor=deepblue,citecolor=deepblue,urlcolor=deepblue,
 pdftitle={Faster Attacks on QR-UOV Signatures},
 pdfauthor={Justin Thaler, Kai Zhang, Scott Kominers, Quang Dao},
 pdfsubject={Algebraic cryptanalysis of QR-UOV},
 pdfkeywords={QR-UOV,UOV,structural cryptanalysis,projective slicing,complete-intersection discriminant,geometric resolution,pseudo-oil}}

\newtheorem{theorem}{Theorem}[section]
\newtheorem{lemma}[theorem]{Lemma}
\newtheorem{proposition}[theorem]{Proposition}
\newtheorem{corollary}[theorem]{Corollary}
\newtheorem{definition}[theorem]{Definition}
\newtheorem{remark}[theorem]{Remark}
\newtheorem{heuristic}[theorem]{Heuristic}

\newcommand{\F}{\mathbb F}
\newcommand{\PP}{\mathbb P}
\newcommand{\cO}{\mathcal O}
\newcommand{\cB}{\mathcal B}
\newcommand{\cX}{\mathcal X}
\newcommand{\cR}{\mathcal R}
\newcommand{\cU}{\mathcal U}
\newcommand{\cW}{\mathcal W}
\newcommand{\cG}{\mathcal G}
\newcommand{\cE}{\mathcal E}
\newcommand{\Span}{\operatorname{span}}
\newcommand{\rank}{\operatorname{rank}}
\newcommand{\codim}{\operatorname{codim}}
\newcommand{\cdeg}{\operatorname{cdeg}}
\newcommand{\Sym}{\operatorname{Sym}}
\newcommand{\disc}{\mathrm{disc}}
\newcommand{\bad}{\mathrm{bad}}
\newcommand{\oil}{\mathrm{oil}}
\newcommand{\off}{\mathrm{off}}
\newcommand{\dep}{\mathrm{dep}}
\newcommand{\QRUOV}{\textsf{QR-UOV}}
\newcommand{\Solve}{\operatorname{Solve}}
\newcommand{\Adv}{\operatorname{Adv}}
\newcommand{\trans}{\mathsf T}
\newcommand{\softO}{\widetilde O}

\setlist[itemize,enumerate]{nosep,leftmargin=*}
\setlength{\parindent}{1.2em}
\setlength{\parskip}{0.15em}
\setlength{\emergencystretch}{2em}
\raggedbottom
\pagestyle{fancy}
\fancyhf{}
\fancyfoot[C]{\thepage}
\renewcommand{\headrulewidth}{0pt}

\title{\bfseries Faster Attacks on \QRUOV{} Signatures}
\author{Justin Thaler\thanks{Georgetown University and a16z crypto research.}
\and Kai Zhang\thanks{MIT.}
\and Scott Kominers\thanks{Harvard University and a16z crypto research.}
\and Quang Dao\thanks{Carnegie Mellon University and LayerZero Labs.}}
\date{24 July 2026}

\begin{document}
\maketitle

\begin{abstract}
We give a faster structural attack on \QRUOV{}, a third-round candidate in NIST's Additional Digital Signature process.  For the Round~2 Levels~I, III, and V parameters, the main projective-slice attack has soft-$O$ bit-complexity exponent indicators
\[
119.59,\qquad 172.71,\qquad 230.17.
\]
A generic $O(T\log T)$ simulation of a $T$-step bit algorithm gives Boolean-circuit exponent indicators
\[
126.49,\qquad 180.14,\qquad 238.02,
\]
against nominal targets $143,207,272$ and the submitters' reported best classical costs $150,211,277$.  These are asymptotic indicators with hidden constants and soft factors, not calibrated gate counts.

The attack restricts the public extension-field quadratic forms to a random projective $V$-plane and solves a square system of $V$ quadrics.  A smooth restricted core contains a unique point on the hidden oil space, and the ambient derivative kernel at that point recovers the full oil space.  We prove that the joint slice-and-output discriminant is nonzero for all but an explicitly bounded fraction of ideal-field keys, below $2^{-364.56},2^{-457.92},2^{-1038.06}$ at the three levels.  A failed public slice is detected and cheaply resampled, so randomization affects the running time rather than correctness.

The recovered object is an \emph{expanded equivalent UOV signing trapdoor}: an oil graph and central matrices that reproduce the public map and permit the UOV signing algebra.  It need not have a preimage under the specification's seed-expansion map, so we do not claim recovery of a valid serialized secret key.  A conditional pseudo-oil direction and a proposed arbitrary-core fallback are separated from the proved main result.
\end{abstract}


\section{Introduction}\label{sec:intro}

Unbalanced Oil and Vinegar (UOV) signatures divide variables into vinegar and oil variables.  In secret coordinates, oil variables never multiply one another.  Equivalently, the public quadratic forms share a hidden linear subspace on which all associated bilinear forms vanish.  Given the corresponding expanded central map, signing chooses vinegar variables and solves a linear system in the oil variables~\cite{KPG1999}.

\QRUOV{} compresses the UOV public key using quotient-ring, or extension-field, structure~\cite{QRUOV2021}.  NIST advanced \QRUOV{} to the third round of its Additional Digital Signature process in May~2026~\cite{NISTRound3,NISTIR8610}.  The Round~2 parameter sets analyzed here are
\[
(q,v,m,\ell)=(127,156,54,3),\quad(127,228,78,3),\quad(127,306,105,3),
\]
for Levels~I, III, and V~\cite{QRUOVSpec}.  As of the manuscript cutoff date, the project website still identifies Specification~v2.0 (5~February~2025) as the current public specification, so these are the latest publicly specified parameter sets analyzed here.  The specification reports best classical attack costs of 150, 211, and 277 bits.

The specification exposes two useful public descriptions.  First, it gives a compact UOV tuple over
\[
K=\F_{127^3},\qquad V=v/3\in\{52,76,102\},\qquad O=m/3\in\{18,26,35\}.
\]
Second, the verifier is a base-field quadratic map with shapes $(n,m)=(210,54),(306,78),(411,105)$.  The proved attack uses the first description.  A secondary, heuristic direct-forgery direction using the second description is recorded separately in \cref{sec:pseudo-oil} and is not used in the security conclusion.

\subsection{Main attack}

The compact public key is a list of $m$ symmetric matrices $P_1,\ldots,P_m$ over $K$, acting on $K^{V+O}$.  The secret structure is an $O$-dimensional subspace $\cO$ on which all $m$ associated bilinear forms vanish.  The attacker first chooses a \emph{regularization field}
\[
L_0=\F_{127^{3s_0}}\supset K
\]
large enough only to make a random evaluation of the joint discriminant succeed with positive, explicitly bounded probability.  It then samples a matrix
\[
M\in L_0^{(V+O)\times(V+1)}.
\]
Its image is a projective $V$-plane.  Since the corresponding vector dimensions are $V+1$ and $O$ inside a space of dimension $V+O$, this plane always meets $\PP(\cO_{L_0})$.

Restricting all $m$ public forms to this plane and taking $V$ random linear combinations gives a square system of $V$ quadrics in $\PP^V_{L_0}$.  When the core is a smooth complete intersection, it consists of exactly $2^V$ reduced geometric points.  Its oil intersection is a single $L_0$-rational projective point: a larger intersection would contain a projective line of common zeros.  At that point the derivatives of all public forms annihilate the oil space, while smoothness gives rank $V$; hence the ambient derivative kernel is exactly $\cO_{L_0}$.  Row reduction recovers the oil graph, Frobenius descent returns it to $K$, and direct substitution verifies every public graph and isotropy identity.

The polynomial-system solve is invoked through Corollary~5.5 of van der Hoeven--Lecerf over the coefficient field $L_0$.  That corollary constructs a further internal extension $\widehat L/L_0$ when needed and retains a bit bound proportional to the input-field degree $[L_0:\F_{127}]=3s_0$.  Because the solver's primitive element and coordinate change may live in $\widehat L$, rational-root extraction is done projectively: after undoing the coordinate change, the attack keeps exactly those roots whose homogeneous coordinate vector is proportional to its $|L_0|$-Frobenius image.  This avoids incorrectly requiring the solver's separator value, or a particular affine representative, to lie in $L_0$.

The attack then reconstructs central UOV matrices and the expanded secret linear map.  This gives the data needed to execute the UOV signing algebra and to produce publicly verifiable signatures.  The official serialized secret key, however, consists of seeds from which the expanded matrix is regenerated.  An arbitrary recovered graph need not be seed-encodable, and even the planted expanded graph does not reveal its seed preimage.  Accordingly, throughout this paper \emph{trapdoor recovery} means recovery of expanded signing data, not recovery of the specification's serialized secret-key bit string.

A fixed public slice does not support a category-negligible exceptional-key theorem: the locus on which its planted derivative is rank deficient has low codimension.  Randomizing both the slice and the $V$ output combinations changes the incidence problem.  We prove that the joint discriminant can vanish identically only in three explicit cases---no usable oil direction, an unwanted off-oil component, or a linear dependence among the public forms.  The dominant case has codimension
\[
e(O+1)+1,\qquad e=m-V,
\]
namely $39,55,109$ for the three levels.  Degree and finite-field point-count bounds turn these codimensions into the bad-key probabilities stated below.

\subsection{Attack at a glance}

\begin{center}
\fbox{\begin{minipage}{0.94\textwidth}
\textbf{Fast expanded-trapdoor recovery.}
\begin{enumerate}
\item Expand the compressed public key to the public extension-field UOV tuple over $K$.  Construct the least extension $L_0/K$ with $|L_0|>3V(V+1)2^{V-1}$.
\item Draw $M\in L_0^{(V+O)\times(V+1)}$ and $R\in L_0^{V\times m}$ once.  A rank failure ends this public trial.  Otherwise restrict all $m$ public quadrics to $\PP(\operatorname{im}M)$ and form the $V$ recombined restricted quadrics.
\item Invoke the certificate-checked small-field Kronecker solver over $L_0$; it may construct an internal field $\widehat L$.  Undo its coordinate change and retain exactly the $L_0$-rational \emph{projective} roots, using Frobenius-proportionality equations on homogeneous coordinates.  Impose all $m$ unrecombined restricted equations.
\item At each survivor, compute the ambient derivative kernel.  Accept only an $O$-space that descends to $K$ and passes all public graph and total-isotropy checks.
\item Reconstruct the expanded central map and secret linear transformation.  Exercise the expanded signing algebra and accept an output only after public signature verification.  No seed preimage or accepted serialized secret key is claimed.
\end{enumerate}
On failure, resample the public coins.  All final acceptance tests are public, so randomization affects termination rather than soundness.
\end{minipage}}
\end{center}

\begin{theorem}[Fast almost-all-key expanded-trapdoor recovery, informal]\label{thm:intro}
For Levels~I, III, and V, choose regularization degrees $[L_0:K]=4,5,6$.  With the single choice $\epsilon=0.2$, the projective-slice attack has soft-$O$ bit-complexity exponent indicators
\[
119.59,\qquad172.71,\qquad230.17,
\]
and generic $O(T\log T)$ circuit-simulation indicators
\[
126.49,\qquad180.14,\qquad238.02.
\]
In the ideal-field experiment, the key-bad probability is below $2^{-364.56},2^{-457.92},2^{-1038.06}$.  For every key outside that set, one public slice-and-output trial fails the smoothness condition with probability at most $2^{-19.85},2^{-15.73},2^{-9.85}$.  Respectively $8,14,28$ independent trials make this algebraic failure smaller than $2^{-143},2^{-207},2^{-272}$; charging all trials in parallel gives circuit indicators $129.49,183.95,242.83$.  Every accepted candidate is checked against the public map, and the planted oil graph is recovered on every smooth trial.
\end{theorem}

\begin{table}[t]
\centering
\small
\setlength{\tabcolsep}{4pt}
\caption{Main projective-slice attack at $\epsilon=0.2$.  ``Direct'' specializes the small-field solver corollary and charges the input coefficient-field degree $3[L_0:K]$.  ``Circuit'' adds the exponent overhead of a generic $O(T\log T)$ simulation.  Neither column is a calibrated operation count.}
\label{tab:headline}
\begin{tabular}{@{}lrrrrrr@{}}
\toprule
Level & $[L_0:K]$ & Direct & Circuit & Target & Spec. best & Margin\\
\midrule
I   & 4 & 119.59 & 126.49 & 143 & 150 & 16.51\\
III & 5 & 172.71 & 180.14 & 207 & 211 & 26.86\\
V   & 6 & 230.17 & 238.02 & 272 & 277 & 33.98\\
\bottomrule
\end{tabular}
\end{table}

\paragraph{Rigor and cost model.}
The projective geometry, ideal-field distribution, incidence bounds, descent, and expanded-trapdoor reconstruction are exact where stated.  Fixed-budget success transfers to the concrete SHAKE/AES-expanded distribution with the distinguishing and rejection terms in \cref{prop:transfer}.  The numerical complexity columns are asymptotic exponent indicators obtained by specializing published bounds; they suppress constants, polylogarithmic factors, memory traffic, and engineering overhead.  The pseudo-oil and arbitrary-core sections are explicitly conditional and do not enter \cref{thm:intro,tab:headline}.

\subsection{Related work}\label{sec:related}

\QRUOV{} was introduced as a quotient-ring compression of UOV at ASIACRYPT~2021~\cite{QRUOV2021}.  The Round~2 specification analyzes direct MQ, reconciliation, intersection, PXL, and rectangular-MinRank attacks~\cite{QRUOVSpec}.  Furue and Ikematsu sharpened the rectangular-MinRank line, and Suzuki et al. extended that framework while reporting that the recommended sets remained above target in their analyses~\cite{FurueIkematsu,Suzuki}.  Our main attack instead exploits the public extension-field tuple, projective slicing, complete polynomial-system solving, and derivative-kernel recovery.

Beullens' intersection attack and subsequent UOV cryptanalysis recover hidden oil geometry by structured linear algebra~\cite{Beullens}.  P\'ebereau showed that one suitable oil vector suffices for polynomial-time expanded-trapdoor recovery~\cite{Pebereau}.  The derivative-kernel theorem used here is a public-coordinate realization of that principle; the one-column completion theorem is a QR-UOV graph-coordinate specialization.

Van der Hoeven and Lecerf give the finite-field Kronecker solver used by the main attack~\cite{vdHL}.  Benoist computes the complete-intersection discriminant degrees used in the regularization analysis~\cite{Benoist}.  The conditional fallback appendix draws on geometric resolutions and symbolic homotopy due to Schost, Safey El Din--Schost, and Vu~\cite{Schost,SafeyElDinSchost,Vu}; the finite-characteristic specialization and its cost ledger in this manuscript are proposed extensions, not ingredients of the main theorem.

May, Ostuzzi, and Ressler's \emph{Just Guess} and Ostuzzi's enriched pseudo-oil framework give improved algorithms for underdetermined MQ~\cite{JustGuess,Ostuzzi}.  Our secondary pseudo-oil section records a QR-UOV specialization and an exact residual reduction conditional on obtaining the required decomposition.  Jin et al. study odd-characteristic UOV families through symmetric algebra~\cite{Jin}.  Side-channel and vinegar-reuse issues are orthogonal and are surveyed by Aulbach, Campos, and Kr\"amer~\cite{Aulbach}.

\section{The public key over the extension field}\label{sec:model}

This section fixes the notation the attack works with. Recall from \cref{sec:intro} that the specification exposes a compact UOV key over the extension field $K=\F_{127^3}$: a list of quadratic forms sharing a hidden oil space. We write that key out and record the recommended parameters. We then set up the probability model used for the key-density bounds---sampling the key material uniformly, rather than from the specification's hash-based expansion---and show in \cref{prop:transfer} how to transfer those bounds back to concrete seeded keys.

The public key-recovery tuple is
\[
P=(p_1,\ldots,p_m),\qquad p_i(x)=x^{\trans}P_ix,
\qquad x\in K^{V+O},
\]
where $P_i$ are symmetric and share a hidden totally isotropic $O$-space $\cO$:
\begin{equation}\label{eq:isotropy}
u^{\trans}P_iv=0\qquad(u,v\in\cO,1\le i\le m).
\end{equation}
The recommended parameters are:
\begin{center}
\begin{tabular}{@{}lrrrrr@{}}
\toprule
Level & $q$ & $\ell$ & $V$ & $O$ & $m$\\
\midrule
I&127&3&52&18&54\\
III&127&3&76&26&78\\
V&127&3&102&35&105\\
\bottomrule
\end{tabular}
\end{center}
The same records appear in the active Round~2 reference implementation~\cite{QRUOVRef}. Throughout,
\[
Q=|K|=127^3=2{,}048{,}383,
\qquad \log_2Q=20.966054\ldots.
\]

\subsection{Ideal-field model and concrete expansion}

\begin{definition}[Consistent ideal-expansion experiment]\label{def:ideal}
Replace the outputs of the distinct domain-separated calls made by \textsf{Expandsk} and \textsf{Expandpk} in one key-generation transcript by independent uniform field elements, sampled once and cached under their call labels.  Every later use of the same public expansion transcript reuses the cached values.  Equivalently, in extension-field notation, sample
\[
G\leftarrow K^{V\times O},\qquad
A_i\leftarrow\Sym_V(K),\qquad
E_i\leftarrow K^{V\times O}\quad(1\le i\le m)
\]
independently and uniformly, set each lower-right public block by the key-generation identity, and expose the resulting expanded public tuple.  The caching clause is essential because key generation and the attacker both expand the same public seed.
\end{definition}

\begin{proposition}[Exact factorization in the ideal-field experiment]\label{prop:ideal-factorization}
In the block notation of \eqref{eq:central}--\eqref{eq:graphid}, the ideal-field experiment samples the secret graph $G$ and the pairs $(A_i,E_i)$ independently and uniformly.  After defining
\[
B_i:=E_i-A_iG,
\]
the random variables $G$ and $(A_i,B_i)_{1\le i\le m}$ are mutually independent and uniform on their respective matrix spaces.  The lower-right public blocks $C_i$ are then determined by \eqref{eq:graphid}.
\end{proposition}

\begin{proof}
Fix $G$ and all $A_i$.  For each $i$, the map $E_i\mapsto B_i=E_i-A_iG$ is a translation of the additive space $K^{V\times O}$ and hence a bijection.  Thus every tuple $(G,(A_i,B_i)_i)$ has exactly one preimage $(G,(A_i,E_i)_i)$ and the same probability.  Since the latter tuple is sampled once from a product of uniform distributions in \cref{def:ideal} and reused consistently, the former tuple is also a product of uniform distributions.  The formula for $C_i$ is the public pull-back identity \eqref{eq:graphid}.
\end{proof}

The specification uses finite SHAKE/AES buffers and rejection sampling. If a call exhausts its buffer it substitutes zero; the buffers are sized so that any single call exhausts with probability at most $2^{-\lambda}$, where $\lambda\in\{128,192,256\}$ is the category security parameter for Levels~I, III, and V. With $2m+1$ such calls, put
\[
\varepsilon_{\rm RS}=(2m+1)2^{-\lambda}.
\]

Here $\Adv^{\rm exp}_{\QRUOV}(D)$ denotes the advantage of a distinguisher $D$ in telling the scheme's complete, consistently reused SHAKE/AES expansion transcript from the cached independent-uniform transcript of \cref{def:ideal}.  Bounding it is a pseudorandomness assumption on the domain-separated expansion function, and every transfer below to concrete seeded keys is conditional on it; the ideal-field bounds themselves are unconditional.

\begin{proposition}[Expansion transfer]\label{prop:transfer}
Let $U$ be any auxiliary public random variable sampled independently of key generation and identically in the concrete and ideal experiments.  For any event $E$ decidable from the expanded key-generation matrices and $U$, there is a distinguisher $D_E$, with the cost of key generation plus the event test, such that
\[
\left|\Pr_{\rm concrete}[E]-\Pr_{\rm ideal}[E]\right|
\le \Adv^{\rm exp}_{\QRUOV}(D_E)+\varepsilon_{\rm RS}.
\]
\end{proposition}

\begin{proof}
Condition on the value of $U$.  Replace the complete concrete expansion transcript by a cached independent-uniform transcript, paying the transcript distinguishing advantage, and abort if any rejection buffer contains too few accepted bytes, paying at most $(2m+1)2^{-\lambda}$.  Conditioned on no abort, the accepted values are independent and uniform across distinct call labels and are reused consistently, giving exactly \cref{def:ideal}.  Averaging over $U$ proves the claim.  In particular, $U$ may contain the public slice, output-mixing, separator, and solver coins of any fixed-budget verified attack experiment.
\end{proof}

The reported rejection-term logarithms are $-121.23,-184.71,-248.28$. We apply \cref{prop:transfer} to the public verified-success experiment, whose test cost is the attack cost, rather than to exact algebraic-locus membership.

\subsection{Central and public graph coordinates}

In secret central coordinates $K^V\oplus K^O$,
\begin{equation}\label{eq:central}
F_i=\begin{pmatrix}A_i&B_i\\B_i^{\trans}&0\end{pmatrix},
\qquad F_i(v,o)=v^{\trans}A_iv+2v^{\trans}B_io,
\end{equation}
with $A_i\in\Sym_V(K)$ and $B_i\in K^{V\times O}$ independent uniform. The planted oil space is $\cO_0=\{0\}\oplus K^O$.

In normalized public coordinates the hidden space is a graph
\begin{equation}\label{eq:graph}
\cO=\left\{\binom{-Gc}{c}:c\in K^O\right\}.
\end{equation}
Writing
\[
P_i=\begin{pmatrix}A_i&E_i\\E_i^{\trans}&C_i\end{pmatrix},
\]
total isotropy is equivalent to
\begin{equation}\label{eq:graphid}
C_i=G^{\trans}E_i+E_i^{\trans}G-G^{\trans}A_iG.
\end{equation}

\begin{proposition}[Verified graph gives an exact expanded UOV decomposition]\label{prop:verified-key}
Let $G'\in K^{V\times O}$ satisfy every identity \eqref{eq:graphid}, and put
\[
B_i'=E_i-A_iG',\qquad
F_i'=\begin{pmatrix}A_i&B_i'\\ {B_i'}^{\trans}&0\end{pmatrix},\qquad
S_{G'}=\begin{pmatrix}I_V&G'\\0&I_O\end{pmatrix}.
\]
Then
\[
P_i=S_{G'}^{\trans}F_i'S_{G'}
\qquad(1\le i\le m).
\]
Consequently, the public pull-back gives an expanded base-field UOV decomposition of the unchanged verification map.  This proposition does not assert that $G'$ has a preimage under \textsf{Expandsk}, or that the official seed-encoded secret-key interface accepts the recovered data.
\end{proposition}

\begin{proof}
Since
\[
S_{G'}^{-1}=\begin{pmatrix}I_V&-G'\\0&I_O\end{pmatrix},
\]
a direct block multiplication gives
\[
S_{G'}^{-\trans}P_iS_{G'}^{-1}
=
\begin{pmatrix}
A_i&E_i-A_iG'\\
E_i^{\trans}-{G'}^{\trans}A_i&
C_i-{G'}^{\trans}E_i-E_i^{\trans}G'+{G'}^{\trans}A_iG'
\end{pmatrix}.
\]
The lower-right block is zero by \eqref{eq:graphid}; the remaining blocks are exactly $F_i'$.  Rearranging proves the displayed factorization.  The specification's public pull-back identifies the structured matrices over $K$ with expanded base-field blocks, so the same identity holds for the base-field verification map.

The official serialized secret key consists of seeds and the signing routine re-expands its secret matrix from those seeds.  A recovered matrix $G'$ therefore cannot in general be passed to that interface.  An attacker can instead execute the same expanded UOV operations directly: choose vinegar values, solve the induced linear system in the oil variables, apply the expanded secret transformation, and verify the resulting signature publicly.  The planted $G$ recovered by the main theorem reproduces the official expanded secret data; for any other verified decomposition, successful signing is treated as an additional public acceptance test rather than assumed.
\end{proof}

\section{Recovery and structural certificates}\label{sec:recovery}

Once the attack has a projective point $x$ lying on the oil space, three things remain: read the full oil space off from $x$, certify and complete it into an expanded UOV decomposition, and confirm that the recovered space is genuinely the hidden one. This section supplies each step---derivative-kernel recovery (\cref{lem:kernel}), one-column graph completion (\cref{thm:completion}), and a uniqueness bound ruling out a second oil space (\cref{thm:uniqueness})---together with a $2^V$ optimality floor for the associated square systems (\cref{thm:bezout}).

For a public common zero $x$, form
\[
C_x=(P_1x\ \cdots\ P_mx)\in L^{(V+O)\times m}.
\]

\begin{lemma}[Derivative-kernel recovery]\label{lem:kernel}
If $x\in\PP(\cO_L)$, then $\cO_L\subseteq\ker C_x^{\trans}$. If $\rank C_x=V$, equality holds.
\end{lemma}

\begin{proof}
For $u\in\cO_L$, \eqref{eq:isotropy} gives $u^{\trans}P_ix=0$ for every $i$. If $\rank C_x=V$, its transpose kernel has dimension $(V+O)-V=O$.
\end{proof}

A smooth restricted square core has derivative rank $V$ at its oil point, while total isotropy bounds the ambient derivative rank by $V$. Thus \cref{lem:kernel} recovers the complete oil space. Row reduction into the public graph chart, Frobenius descent to $K$, verification of \eqref{eq:graphid}, and \cref{prop:verified-key} recover the expanded trapdoor.  A trial signature is then accepted only after public verification.

\subsection{One-column completion}

Let $G=[g_1|\cdots|g_O]$. For a coordinate-slice candidate $g$, define
\[
(M_j(g))_{i,*}=(A_ig-E_ie_j)^{\trans},
\qquad
(H_j(g))_{i,r}=(E_ie_r)^{\trans}g-(C_i)_{jr}.
\]

\begin{theorem}[One column completes the graph]\label{thm:completion}
For the true column $g_j$,
\[
M_j(g_j)G_{-j}=H_j(g_j).
\]
If $\rank M_j(g_j)=V$, all remaining columns are uniquely recovered by one public multi-right-hand-side solve.
\end{theorem}

\begin{proof}
The $(j,r)$ entry of \eqref{eq:graphid} rearranges to
\[
(A_ig_j-E_ie_j)^{\trans}g_r=(E_ie_r)^{\trans}g_j-(C_i)_{jr}.
\]
Stacking over $i,r$ gives the matrix identity. Full column rank gives uniqueness.
\end{proof}

\subsection{Generic uniqueness}

We now check that the space the attack recovers is the true oil space: for almost every key there is no second common isotropic $O$-space that could be confused with it. Fix the planted oil space $\cO_0$.  Let $W\ne\cO_0$ be another $O$-space and put
\[
s=O-\dim(W\cap\cO_0),\qquad 1\le s\le O.
\]

\begin{lemma}[Relative Grassmannian stratum]\label{lem:relative-grassmann}
The locus of such $W$ with fixed $s$ has dimension $s(V+O-s)$.
\end{lemma}

\begin{proof}
Choose the $(O-s)$-dimensional intersection inside $\cO_0$, the $s$-dimensional image of $W$ in $K^{V+O}/\cO_0\simeq K^V$, and the graph map back to the $s$-dimensional quotient of $\cO_0$.  The dimensions are $(O-s)s$, $s(V-s)$, and $s^2$, respectively.
\end{proof}

\begin{lemma}[Independent extra isotropy conditions]\label{lem:extra-isotropy}
For fixed $W$ in this stratum, requiring one central form that already vanishes on $\cO_0\times\cO_0$ also to vanish on $W\times W$ imposes exactly
\[
c_s=\binom{O+1}{2}-\binom{O-s+1}{2}=\frac{s(2O-s+1)}2
\]
additional independent linear conditions.
\end{lemma}

\begin{proof}
Restriction to $W$ already vanishes on $\Sym^2(W\cap\cO_0)$.  Conversely, after choosing complements, every symmetric form on $W$ that vanishes on this subspace extends to an ambient symmetric form vanishing on $\cO_0\times\cO_0$.  The restriction map therefore has the displayed rank.
\end{proof}

\begin{theorem}[Uniqueness codimension and rational first moment]\label{thm:uniqueness}
For $m$ independent central forms, the key locus admitting an alternative common isotropic $O$-space of type $s$ has codimension at least
\[
\delta_s=m\frac{s(2O-s+1)}2-s(V+O-s).
\]
The minimum over $s$ equals $903,1927,3539$ for Levels~I, III, and V.  Moreover, with Gaussian binomial coefficients denoted by $\genfrac{[}{]}{0pt}{}{a}{b}_Q$,
\[
\Pr[\text{a second $K$-rational oil space}]
\le
\sum_{s=1}^{O}
\genfrac{[}{]}{0pt}{}{O}{s}_Q
\genfrac{[}{]}{0pt}{}{V}{s}_Q
Q^{s^2-mc_s}.
\]
The base-two logarithms of this bound are $-18932.35,-40401.59,-74198.87$.
\end{theorem}

\begin{proof}
Over the stratum of dimension $s(V+O-s)$, each of the $m$ form fibers loses $c_s$ dimensions by \cref{lem:extra-isotropy}.  The incidence variety therefore has codimension at least $mc_s-s(V+O-s)$ after projection to key space.  The function $\delta_s$ is concave in $s$, so its minimum on $1\le s\le O$ occurs at an endpoint; direct evaluation gives the displayed values and shows that $s=1$ minimizes in all three parameter sets.

For the finite-field statement, the number of $K$-rational $O$-spaces in the type-$s$ stratum is
\[
\genfrac{[}{]}{0pt}{}{O}{s}_Q
\genfrac{[}{]}{0pt}{}{V}{s}_Q Q^{s^2}.
\]
For each fixed $W$, the $m$ independent forms satisfy the $c_s$ extra conditions with probability $Q^{-mc_s}$.  The union bound gives the formula; exact evaluation gives the stated logarithms.
\end{proof}

\subsection{Multihomogeneous path-count floor}

For a square subsystem of the coupled graph equations, let $d_j$ count selected diagonal equations in block $j$ and $c_{jk}$ selected cross equations. Its multihomogeneous B\'ezout number is
\[
B=[t_1^V\cdots t_O^V]\prod_j(2t_j)^{d_j}\prod_{j<k}(t_j+t_k)^{c_{jk}}.
\]

\begin{theorem}[B\'ezout floor]\label{thm:bezout}
If $B\ne0$, then $B\ge2^V$. Equality is attained by taking $V$ diagonal equations in one root block and $V$ cross equations on each edge of a rooted spanning tree. Thus one-column recovery followed by linear completion is optimal in multihomogeneous path count among square graph subsystems.
\end{theorem}

\begin{proof}
Expand each cross factor by orienting a labeled edge toward the endpoint whose variable is selected. If $D=\sum_jd_j$, then $B=2^DN$, where $N$ counts orientations with indegrees $V-d_j$. If $D\ge V$, the claim is immediate. Otherwise every vertex has indegree at least $V-D$.  Remove a directed cycle; after fewer than $V-D$ such removals, each vertex has lost at most one incoming edge per removed cycle and still has positive indegree, so another directed cycle remains.  Hence there are $V-D$ edge-disjoint directed cycles.  Reversing any subset preserves all indegrees and gives $N\ge2^{V-D}$. The rooted-tree construction has one orientation and $D=V$.
\end{proof}

\section{Random projective slicing and fast regularization}\label{sec:projective}

This section proves the regularization and key-density claims behind \cref{thm:intro}. We work in secret central coordinates, where the key distribution is transparent, and show that for all but a negligible fraction of keys a random slice together with a random recombination produces a smooth $2^V$-point core from which the oil space is recovered.

Let $K=\F_{127^3}$, $Q=|K|$, and
\[
(V,O,m)\in\{(52,18,54),(76,26,78),(102,35,105)\}.
\]
Put $e=m-V\in\{2,2,3\}$ and $N=V+O$.  Work first in secret central coordinates
\[
W=K^V\oplus K^O,
\qquad
\cO_0=\{0\}\oplus K^O.
\]
The $m$ central quadrics are
\begin{equation}\label{eq:central-slice}
F_i(v,o)=v^{\mathsf T}A_i v+2v^{\mathsf T}B_i o,
\qquad
A_i\in\Sym_V(K),\ B_i\in K^{V\times O}.
\end{equation}
In the ideal-field experiment, the pairs $(A_i,B_i)$ are independent and uniform.  Equivalently, each $F_i$ is uniform in
\[
\cU:=H^0\!\left(\PP(W),\mathcal I_{\PP(\cO_0)}(2)\right),
\qquad
D_{\rm form}:=\dim\cU=\frac{V(V+1)}2+VO.
\]
The proof is written in central coordinates only to expose the distribution.  The secret change from public to central coordinates lies in $\operatorname{GL}_N(K)$; left multiplication by its inverse is a bijection on $L_0^{N\times(V+1)}$.  Thus a public uniform slice matrix remains uniform in central coordinates, while smoothness, intersection multiplicity, and derivative-kernel recovery are coordinate invariant.  In particular, the public and central tuples have identically vanishing joint discriminant polynomials simultaneously.
Let
\[
X_F:=V(F_1,\ldots,F_m)\subset\PP(W).
\]
It contains the planted oil space $\PP(\cO_0)\simeq\PP^{O-1}$.

The fast branch samples
\[
M\leftarrow L_0^{N\times(V+1)}
\]
and declares the public trial failed if $\rank M<V+1$, where $L_0/K$ is the regularization field chosen below.  It otherwise restricts the public quadrics to the projective $V$-plane $\Lambda_M=\PP(\operatorname{im}M)$:
\[
q_i(t)=F_i(Mt),\qquad t\in\PP^V_{L_0},
\]
and then samples $R\leftarrow L_0^{V\times m}$, declares failure if $\rank R<V$, and solves
\[
g_a(t)=\sum_{i=1}^m R_{a,i}q_i(t),\qquad 1\le a\le V.
\]
Treating a rank defect as a failed draw, rather than conditioning the distribution on full rank, lets the unconditioned Schwartz--Zippel bound apply verbatim; by \cref{lem:theta-ranks}, every such defect already lies in $\{\Theta_F=0\}$.

In secret affine coordinates, a transverse $\Lambda_M$ is exactly a slice $o=c+Lv$ after translating its unique oil intersection.  The projective parameterization is preferable in the proof because it has no inverses or denominators: every restricted coefficient is simply quadratic in the entries of $M$.

\subsection{Why the full discriminant is enough}

The van der Hoeven--Lecerf solver is stated under regular-sequence and absolute-radicality assumptions for all intermediate ideals.  It is not necessary to force these assumptions by multiplying all prefix discriminants.

\begin{lemma}[Full smooth core implies the solver's prefix hypotheses]\label{lem:full-implies-prefix}
Let $g_1,\ldots,g_V$ be quadrics in $\PP^V$ over a field of odd characteristic.  If their common zero scheme consists of exactly $2^V$ distinct points and the projective Jacobian has rank $V$ at every point, then the ordered system satisfies the regular-sequence and absolutely-radical intermediate-ideal hypotheses used by the Kronecker solver.
\end{lemma}

\begin{proof}
Extend scalars to an algebraic closure.  Choose a projective coordinate change whose affine chart contains all $2^V$ points, and choose a linear form separating them.  Interpolation then gives a separable polynomial $q$ of degree $2^V$ and coordinate polynomials $v_1,\ldots,v_V$ parametrizing all roots.  The original equations vanish modulo $q$, and the Jacobian determinant is invertible modulo $q$ because every root is regular.  These are exactly the hypotheses of the a posteriori criterion of van der Hoeven and Lecerf~\cite[Proposition~4.6]{vdHL}; it concludes that the sequence is regular and every intermediate ideal is absolutely radical.  These properties are intrinsic and hence descend to the original field.  The solver's remaining degree condition holds because every equation is quadratic.
\end{proof}

Let $\Delta_{V,V}$ be the reduced discriminant for $V$ quadrics in $\PP^V$.  It vanishes exactly when the common zero scheme is not a smooth complete intersection of codimension $V$.  Benoist's degree formula~\cite{Benoist} gives the partial degree
\begin{equation}\label{eq:dpartial}
d_V=(V+1)2^{V-1}
\end{equation}
in the coefficients of each equation.  Thus it suffices to make the single polynomial
\begin{equation}\label{eq:theta}
\Theta_F(M,R):=
\Delta_{V,V}\!\left(
\sum_iR_{1i}F_i(Mt),\ldots,
\sum_iR_{Vi}F_i(Mt)
\right)
\end{equation}
nonzero.

\begin{lemma}[Nonvanishing certifies full parameter rank]\label{lem:theta-ranks}
If $\Theta_F(M,R)\ne0$, then $\rank M=V+1$ and $\rank R=V$.
\end{lemma}

\begin{proof}
If $\rank M<V+1$, then $\PP(\ker M)$ is nonempty and every pull-back $F_i(Mt)$, together with all of its first derivatives, vanishes there; the core is singular.  If $\rank R<V$, the output equations are linearly dependent and cannot cut a smooth codimension-$V$ complete intersection.  In either case the full discriminant vanishes.
\end{proof}

\subsection{Geometric nonvanishing of the joint slice-and-output discriminant}

For $x\in X_F$, let
\[
\rho_F(x):=\rank\bigl(dF_1(x),\ldots,dF_m(x)\bigr),
\]
where the differentials are taken on $T_x\PP(W)$.

\begin{definition}[Three key-level bad loci]\label{def:bad-loci}
Define:
\begin{align*}
\cB_{\oil}&:=\{F:\rho_F(p)<V\text{ for every }p\in\PP(\cO_0)(\overline K)\},\\
\cB_{\off}&:=\{F:\exists x\in X_F(\overline K)\setminus\PP(\cO_0)\text{ with }\rho_F(x)<V\},\\
\cB_{\dep}&:=\{F:F_1,\ldots,F_m\text{ are linearly dependent in }\cU\}.
\end{align*}
\end{definition}

The next three lemmas isolate the only geometry needed.

\begin{lemma}[A transverse finite base slice through a regular oil point]\label{lem:linear-section}
Assume $F\notin\cB_{\oil}\cup\cB_{\off}$.  Choose
$p\in\PP(\cO_0)(\overline K)$ with $\rho_F(p)=V$.  Then a Zariski-open dense set of projective $V$-planes $\Lambda\subset\PP(W_{\overline K})$ through $p$ has all of the following properties:
\begin{enumerate}
\item $\Lambda\cap\PP(\cO_0)=\{p\}$;
\item $X_F\cap\Lambda$ is finite;
\item at every $x\in X_F\cap\Lambda$, the restricted differential $T_x\Lambda\to\overline K^m$ has rank $V$.
\end{enumerate}
\end{lemma}

\begin{proof}
First, $\dim X_F\le O-1$.  Otherwise a component of dimension at least $O$ has a smooth point $x$ outside $\PP(\cO_0)$, and then
\[
\rho_F(x)\le (N-1)-\dim T_xX_F\le V-1,
\]
contrary to $F\notin\cB_{\off}$.  Since $\PP(\cO_0)\subset X_F$, the maximal dimension is exactly $O-1$.

At $p$, every differential annihilates $T_p\PP(\cO_0)$.  Both this tangent space and $\ker dF(p)$ have dimension $O-1$, so they are equal.  Therefore a general $V$-plane through $p$ is transverse at $p$ and intersects the oil space only in $p$.  We henceforth restrict to this nonempty open subset of the Grassmannian; no other oil point can occur in the slice.

Let $\mathcal G_p$ be the Grassmannian of projective $V$-planes through $p$, of dimension $V(O-1)$.  We now consider only $x\in X_F\setminus\PP(\cO_0)$.  Put $K_x=\ker dF(x)$; the off-oil hypothesis gives $\dim K_x\le O-1$.  For fixed $x$, planes through $p$ and $x$ form a family of dimension $(V-1)(O-1)$.  If the tangent direction of the line $\overline{px}$ is not in $K_x$, the additional condition $T_x\Lambda\cap K_x\ne0$ is a proper Schubert condition and lowers this dimension by at least one.  Allowing $x$ to vary in dimension at most $O-1$ gives total dimension at most $V(O-1)-1$.

If the tangent direction of $\overline{px}$ lies in $K_x$, then every quadric $F_i$ vanishes identically on that line: both endpoint values and the polar cross term vanish.  The variety of such line directions through $p$ has dimension at most $\dim X_F-1\le O-2$.  Planes containing one such line therefore again form a family of dimension at most
\[
(O-2)+(V-1)(O-1)=V(O-1)-1.
\]
The off-oil rank-defect incidence is locally closed and has dimension at most $V(O-1)-1$ by the two cases above.  Its closure in the projective product $X_F\times\mathcal G_p$ has the same dimension, and the projection of that closure is therefore a proper closed subset of $\mathcal G_p$.  Outside it, every intersection point has restricted derivative rank $V$.  A positive-dimensional intersection would have a smooth off-oil point with a nonzero tangent direction in $T_x\Lambda\cap\ker dF(x)$, a contradiction.  Hence the intersection is finite.
\end{proof}

\begin{lemma}[A second-order spare equation forces separability]\label{lem:second-order}
Let $f_1,\ldots,f_m$ be quadrics on $\PP^V$ with a common point $p$.  Suppose their differential matrix at $p$ has rank $V$.  If there is a nonzero
\[
h\in\Span(f_1,\ldots,f_m)
\]
whose first jet at $p$ is zero, then the rational map
\[
\phi_f:\PP^V\dashrightarrow\PP^{m-1},
\qquad
x\longmapsto[f_1(x):\cdots:f_m(x)]
\]
has differential rank $V$ on a nonempty open subset.
\end{lemma}

\begin{proof}
Choose projective coordinates with $p=[1:0:\cdots:0]$ and affine coordinates $y=(y_1,\ldots,y_V)$.  After an equation-basis change, the first $V$ forms have expansions
\[
f_j(y)=y_j+q_j(y),
\]
where each $q_j$ is quadratic.  Since $h$ is a quadric with zero value and zero differential at $p$, it is a nonzero homogeneous quadratic $q(y)$ on this chart.

Consider the $(V+1)\times(V+1)$ first-jet minor whose rows are $f_1,\ldots,f_V,h$ and whose columns are value followed by the $V$ first derivatives.  Up to a nonzero sign, its lowest-degree homogeneous term is
\[
q(y)-\nabla q(y)\cdot y=-q(y),
\]
using Euler's identity $\nabla q(y)\cdot y=2q(y)$.  It is therefore nonzero.  Hence the full first-jet rank is $V+1$ generically, which is equivalent to differential rank $V$ for $\phi_f$ away from the base locus.
\end{proof}

\begin{lemma}[Finite regular base locus plus separability]\label{lem:finite-base-mixing}
Let $f_1,\ldots,f_m$ be quadrics on $\PP^V$ over an algebraically closed field.  Assume:
\begin{enumerate}
\item their common base scheme $B=V(f_1,\ldots,f_m)$ is finite and nonempty;
\item the $m\times V$ projective differential matrix has rank $V$ at every point of $B$;
\item $\phi_f$ has differential rank $V$ on a nonempty open subset.
\end{enumerate}
Then a nonempty Zariski-open set of matrices $R\in k^{V\times m}$ makes
\[
\left(\sum_iR_{1i}f_i,\ldots,\sum_iR_{Vi}f_i\right)
\]
a smooth zero-dimensional complete intersection of degree $2^V$.
\end{lemma}

\begin{proof}
Let $\cR=k^{V\times m}$.  For $x\notin B$, write
\[
a(x)=(f_1(x),\ldots,f_m(x))\ne0
\]
and let $J(x)$ be the differential matrix.  A matrix $R$ has $x$ as a root exactly when $Ra(x)=0$, a codimension-$V$ linear condition on $\cR$.

On the open set where the first-jet matrix $[a(x)\mid J(x)]$ has rank $V+1$, the map
\[
a(x)^\perp\longrightarrow T_x^*\PP^V,
\qquad
r\longmapsto rJ(x)
\]
is surjective.  Consequently the additional condition $\det(RJ(x))=0$ has codimension one inside $Ra(x)=0$.  After allowing $x$ to vary in a $V$-dimensional space, this singular-root incidence still has codimension at least one in $\cR$.

The first-jet rank-defect locus away from $B$ has dimension at most $V-1$ by assumption~3.  There the root condition $Ra(x)=0$ alone has codimension $V$, so its total incidence also has dimension at most $\dim\cR-1$.

Finally, at each $x\in B$, the root condition is automatic, but the rank-$V$ derivative assumption makes $\det(RJ(x))=0$ a genuine hypersurface in $\cR$.  Since $B$ is finite, these basepoint hypersurfaces do not fill $\cR$.

Globally, the singular-root incidence is the closed subset of $\PP^V\times\cR$ cut out by $Ra(x)=0$ and the maximal minors of $RJ(x)$.  Its projection is closed because $\PP^V$ is proper, and the preceding stratified dimension bounds show that this projection is proper.  Choose $R$ outside it.  The resulting projective zero scheme is nonempty because it contains $B$.  It cannot have positive dimension: a positive-dimensional component has a smooth point at which the Jacobian rank is below $V$.  Therefore the zero scheme is smooth and zero-dimensional.  Since it is cut out by $V$ quadrics in $\PP^V$, it is a complete intersection of degree $2^V$.
\end{proof}

\begin{theorem}[The enlarged exceptional locus is explicitly controlled]\label{thm:joint-nonzero}
If
\[
F\notin\cB_{\oil}\cup\cB_{\off}\cup\cB_{\dep},
\]
then the joint polynomial $\Theta_F(M,R)$ in \eqref{eq:theta} is nonzero.
\end{theorem}

\begin{proof}
Choose a regular oil point $p$.  By \cref{lem:linear-section}, a dense open subset of the irreducible Grassmannian of $V$-planes through $p$ has finite base intersection and full restricted derivative rank at every basepoint.

The coefficient kernel of the derivative map at $p$ has dimension $m-V=e>0$.  Because $F\notin\cB_{\dep}$, any nonzero vector in this kernel produces a nonzero global quadric $H\in\Span(F_i)$ with zero first jet at $p$.  The condition $H|_\Lambda\not\equiv0$ is open and nonempty: if $H$ vanished on every $V$-plane through $p$, it would vanish on their union, which is all of $\PP(W)$.  Intersecting the two nonempty open subsets gives a plane $\Lambda$ satisfying both properties.

On this plane, \cref{lem:second-order} makes the restricted linear system generically separable, and \cref{lem:finite-base-mixing} supplies an output matrix $R$ whose full core is a smooth complete intersection of degree $2^V$.  Thus the discriminant evaluates nontrivially at one pair $(M,R)$, proving that $\Theta_F$ is not the zero polynomial.
\end{proof}

\subsection{Finite-field density of the three bad loci}

\subsubsection{First-jet surjectivity away from oil}

\begin{lemma}[The central-form space spans every off-oil first jet]\label{lem:jet-surj}
For every $x\in\PP(W)\setminus\PP(\cO_0)$, the first-jet evaluation map
\[
\cU\longrightarrow K\oplus T_x^*\PP(W)
\]
is surjective after scalar extension to $\overline K$.
\end{lemma}

\begin{proof}
Choose a vinegar coordinate $v_j$ that is nonzero at $x$ and normalize to the chart $v_j=1$.  The quadrics
\[
v_j^2,\quad v_jv_k\ (k\ne j),\quad v_jo_a\ (1\le a\le O)
\]
belong to $\cU$ and restrict to the affine functions $1$, the remaining vinegar coordinates, and the oil coordinates.  Their first jets form a basis of the value-and-cotangent space.
\end{proof}

\subsubsection{The off-oil incidence}

Let $\mathcal A=\cU^m$ be the affine key space.  For an affine variety, write $\cdeg$ for the sum of the degrees of its irreducible components.  For each vinegar chart $C_j\simeq\mathbb A^{N-1}$, let $Z_j\subset\mathcal A\times C_j$ be the incidence of pairs $(F,x)$ such that all $F_i(x)$ vanish and the $m\times(N-1)$ derivative matrix has rank at most $V-1$.

\begin{lemma}[Codimension and cumulative degree of the off-oil incidence]\label{lem:off-incidence}
Each $Z_j$ is pure of codimension
\[
m+c_{\det},\qquad c_{\det}=(m-V+1)O=(e+1)O,
\]
and satisfies
\[
\cdeg(Z_j)\le 3^m(2V)^{c_{\det}}.
\]
If $Y_j$ is the Zariski closure of its projection to key space, then
\[
\codim_{\mathcal A}Y_j\ge c_{\off}:=m+c_{\det}-(N-1)=e(O+1)+1,
\qquad
\cdeg(Y_j)\le\cdeg(Z_j).
\]
\end{lemma}

\begin{proof}
On the chart $v_j=1$, the fixed monomials $v_j^2$, $v_jv_k$ for $k\ne j$, and $v_jo_a$ restrict to the affine functions $1$, the remaining vinegar coordinates, and the oil coordinates.  They therefore give a polynomial splitting of the value-and-first-jet evaluation map.  After a triangular polynomial change of coefficient coordinates, the incidence is the product of an affine space, the equations setting $m$ value coordinates to zero, and the determinantal variety of $m\times(N-1)$ matrices of rank at most $V-1$.  It is consequently pure of the stated codimension.

In the original joint coefficient and point coordinates, each value equation has total degree at most three, each derivative entry has degree at most two, and every $V\times V$ minor has degree at most $2V$.  The splitting above shows that $Z_j$ is irreducible and that the value equations and determinantal conditions have the asserted total codimension.  Over the algebraic closure, choose $c_{\det}$ generic linear combinations of the maximal minors so that, after the $m$ value equations, they cut properly and retain $Z_j$ as an irreducible component.  Cumulative B\'ezout therefore gives
\[
\cdeg(Z_j)\le 3^m(2V)^{c_{\det}}
\]
by the cumulative-degree theorem of Heintz and Lachaud--Rolland~\cite{Heintz,LachaudRolland}.

Projection cannot increase dimension, so the projected codimension is at least the source codimension minus the chart dimension $N-1$.  For the degree statement, let $Y$ be an irreducible component of $Y_j$ and choose an irreducible component $Z$ of $Z_j$ that dominates it.  Intersect $Z$ with $\dim Z-\dim Y$ general affine hyperplanes so that the resulting $Z'$ still dominates $Y$ generically finitely and $\deg Z'\le\deg Z$.  A general complementary linear section of $Y$ pulls back to at least $\deg Y$ points of $Z'$, counted with multiplicity, while their number is at most $\deg Z'$.  Thus $\deg Y\le\deg Z$; summing over image components gives $\cdeg(Y_j)\le\cdeg(Z_j)$.
\end{proof}

\begin{proposition}[Explicit off-oil density]\label{prop:off-density}
Under the ideal-field distribution,
\begin{equation}\label{eq:etaoff}
\Pr[F\in\cB_{\off}]
\le
\eta_{\off}:=
V\,3^m(2V)^{(e+1)O}
Q^{-\{e(O+1)+1\}}.
\end{equation}
\end{proposition}

\begin{proof}
Every off-oil point lies in one of the $V$ vinegar charts, so $\cB_{\off}\subseteq\bigcup_jY_j$.  The affine finite-field point bound
\[
|Y_j(\F_Q)|\le\cdeg(Y_j)Q^{\dim Y_j}
\]
from Lachaud--Rolland~\cite{LachaudRolland}, followed by the union bound and \cref{lem:off-incidence}, gives the formula after division by $|\mathcal A(\F_Q)|$.
\end{proof}

\subsubsection{The oil and dependence loci}

Let
\[
\tau_{m,V}=1-\prod_{j=0}^{V-1}(1-Q^{j-m})
\]
be the rank-deficiency probability of a uniform $m\times V$ matrix.  If every geometric oil direction is rank deficient, then in particular all $O$ standard oil-coordinate directions are rank deficient.  Their derivative matrices are independent in the ideal-field experiment.  Hence
\begin{equation}\label{eq:etaoil}
\Pr[F\in\cB_{\oil}]\le\eta_{\oil}:=\tau_{m,V}^{O}.
\end{equation}

The $m$ central forms are independent uniform vectors in a $D_{\rm form}$-dimensional space.  Hence
\begin{equation}\label{eq:etadep}
\Pr[F\in\cB_{\dep}]
=1-\prod_{j=0}^{m-1}(1-Q^{j-D_{\rm form}})
\le
\eta_{\dep}:=
Q^{-D_{\rm form}}\frac{Q^m-1}{Q-1}.
\end{equation}

Combining \cref{thm:joint-nonzero} with
\eqref{eq:etaoff}, \eqref{eq:etaoil}, and \eqref{eq:etadep} gives the promised key theorem.

\begin{theorem}[Almost-all-key joint regularizability]\label{thm:key-density}
In the ideal-field experiment,
\begin{equation}\label{eq:keybad-total}
\Pr[\Theta_F\equiv0]
\le\eta_{\off}+\eta_{\oil}+\eta_{\dep}.
\end{equation}
For the recommended parameters:
\begin{center}
\begin{tabular}{@{}lrrrr@{}}
\toprule
Level & $c_{\off}$ & $\log_2\eta_{\off}$ & $\log_2\eta_{\oil}$ & $\log_2\Pr[\Theta_F\equiv0]$\\
\midrule
I   & 39  & $-364.564$  & $-1132.167$ & $<-364.563$\\
III & 55  & $-457.920$  & $-1635.352$ & $<-457.919$\\
V   & 109 & $-1038.067$ & $-2935.248$ & $<-1038.066$\\
\bottomrule
\end{tabular}
\end{center}
The dependence term is below $2^{-4.7\times10^4}$ already at Level I and is invisible at the displayed precision.
\end{theorem}

\begin{proof}
By \cref{thm:joint-nonzero}, the event $\Theta_F\equiv0$ is contained in $\cB_{\oil}\cup\cB_{\off}\cup\cB_{\dep}$.  Apply \eqref{eq:etaoff}, \eqref{eq:etaoil}, \eqref{eq:etadep}, and the union bound.  The displayed values are exact evaluations of those formulas.
\end{proof}

\subsection{Regularization field, solver extension, and one-trial probability}

Let $d_V$ be as in \eqref{eq:dpartial}.  A coefficient of $F_i(Mt)$ is quadratic in $M$, while the $a$th combined equation is linear in the $a$th row of $R$.  Since $\Delta_{V,V}$ has degree $d_V$ in each equation, the joint polynomial obeys
\begin{equation}\label{eq:joint-degrees}
\deg_M\Theta_F\le2Vd_V,
\qquad
\deg_R\Theta_F\le Vd_V,
\qquad
\deg\Theta_F\le3V(V+1)2^{V-1}.
\end{equation}

Put
\[
D_\Theta(V):=3V(V+1)2^{V-1}
\]
and choose the least $s_0\ge1$ for which $Q^{s_0}>D_\Theta(V)$.  The \emph{regularization field} is
\[
L_0=\F_{Q^{s_0}}=\F_{127^{3s_0}}.
\]
It is used for the public slice and output-recombination coins, and hence is also the coefficient field of the restricted square system.  The selected degrees are much smaller than those obtained by forcing the coefficient field itself to meet the large-field hypothesis of the solver's Theorem~5.4:
\begin{center}
\begin{tabular}{@{}lrrr@{}}
\toprule
Level & $s_0=[L_0:K]$ & $\log_2D_\Theta(V)$ & $\log_2|L_0|$\\
\midrule
I   & 4 & 64.013  & 83.864\\
III & 5 & 89.100  & 104.830\\
V   & 6 & 115.944 & 125.796\\
\bottomrule
\end{tabular}
\end{center}

\begin{corollary}[Joint slice-and-output failure]\label{cor:joint-sz}
For every key outside the bad set of \cref{thm:key-density}, if $M$ and $R$ are sampled uniformly once over $L_0$, with a rank defect counted as failure, then
\begin{equation}\label{eq:delta}
\Pr[\Theta_F(M,R)=0]
\le
\delta_V(L_0):=
\frac{3V(V+1)2^{V-1}}{|L_0|}.
\end{equation}
The numerical bounds and amplification counts are:
\begin{center}
\begin{tabular}{@{}lrrrrr@{}}
\toprule
Level & $[L_0:K]$ & $\log_2\delta_V(L_0)$ & $(1-\delta_V)^{-1}$ & $t_{\rm target}$ & $t_{\rm RS}$\\
\midrule
I   & 4 & $-19.851$ & $1.000001$ & 8  & 7\\
III & 5 & $-15.731$ & $1.000018$ & 14 & 12\\
V   & 6 & $-9.852$  & $1.001083$ & 28 & 26\\
\bottomrule
\end{tabular}
\end{center}
Here $t_{\rm target}$ is the least number of independent public trials for which $\delta_V^{t}$ is below the nominal category target, and $t_{\rm RS}$ is the least number for which it is below the rejection-transfer term $\varepsilon_{\rm RS}$.
\end{corollary}

\begin{proof}
For a nonexceptional key, $\Theta_F$ is a nonzero polynomial by \cref{thm:key-density}.  Its total degree is bounded by \eqref{eq:joint-degrees}, so Schwartz--Zippel gives \eqref{eq:delta}.  Rank-deficient $M$ or $R$ already lie in its zero set by \cref{lem:theta-ranks}; no conditioning or separate rank-rejection term is needed.  The remaining entries are exact evaluations of $\delta_V$, a geometric-series expectation, and the indicated ceiling formulas.
\end{proof}

\paragraph{Small-field solver interface.}
After $M$ and $R$ have been chosen, the restricted quadrics have coefficients in $L_0$.  We invoke Corollary~5.5 of van der Hoeven--Lecerf with input field $L_0=\F_{127^{3s_0}}$~\cite{vdHL}.  Under the regular-sequence, absolute-radicality, and degree hypotheses, that corollary constructs an integer $\ell_{\rm int}=O(V)$, an internal finite extension
\[
\widehat L=\F_{127^{3s_0\ell_{\rm int}}}\supset L_0,
\]
a coordinate change and primitive element over $\widehat L$, and a Kronecker parametrization of the geometric zero set.  Its final Las Vegas bit bound is
\[
\softO\!\left(D\,2^{(1+\epsilon)(V-1)}(3s_0)\log 127\right),
\qquad D=2^V,
\]
for every fixed rational $\epsilon>0$.  Thus the internal extension construction is part of the cited bound; one must not multiply the final indicator by its degree a second time.

A smooth full core has all $2^V$ roots distinct and regular, so \cref{lem:full-implies-prefix} supplies the solver's regular-sequence and absolute-radicality hypotheses.  The solver's coordinate changes, separating forms, and certification coins are independent of the public regularization variables $(M,R)$; failed certificates trigger retry.  The key-bad term from \cref{thm:key-density} is much smaller than every displayed public-trial term and does not change the numerical indicators.

\subsubsection{Field-aware fast-solver cost specialization}

For quadrics the B\'ezout number is $D=2^V$.  Specializing Corollary~5.5 over the input field $L_0=\F_{127^{3s_0}}$ gives the base-two bit-complexity exponent indicator
\begin{equation}\label{eq:fast-cost}
C^{\rm bit}_{\epsilon}(V,s_0)
=V+(1+\epsilon)(V-1)+\log_2\!\bigl(3s_0\log_2 127\bigr).
\end{equation}
The factor $s_0$ is essential because the restricted coefficients are arbitrary elements of $L_0$, while the solver's further extension is already covered by the small-field corollary.  We use the single predeclared value $\epsilon=0.2$ for all three levels.  The direct indicators are
\[
119.59,\qquad172.71,\qquad230.17.
\]
A generic $O(T\log T)$ simulation gives the circuit-size exponent indicator
\begin{equation}\label{eq:circuit-cost}
C^{\rm circ}_{\epsilon}(V,s_0)
=C^{\rm bit}_{\epsilon}(V,s_0)+\log_2 C^{\rm bit}_{\epsilon}(V,s_0),
\end{equation}
namely
\[
126.49,\qquad180.14,\qquad238.02.
\]
These formulas are asymptotic upper-bound indicators.  Hidden constants and soft factors prevent interpreting them as literal calibrated gate counts.  Truncating the Las Vegas solver after a constant multiple of its expected time gives a fixed-budget constant-success computation without changing the exponent.  The expected public regularization retry factors are displayed in \cref{cor:joint-sz} and are at most $1.0011$.

For a bounded public experiment whose algebraic failure is below the nominal target, run $t_{\rm target}=8,14,28$ trials in parallel.  Charging the full factor $t_{\rm target}$ adds $\log_2t_{\rm target}$ to \eqref{eq:circuit-cost}, giving
\[
129.49,\qquad183.95,\qquad242.83,
\]
still below the targets by $13.51,23.05,29.17$ bits under the same indicator convention.  As a sensitivity check, the unamplified circuit indicators remain below target for every fixed $\epsilon$ below approximately $0.520,0.555,0.534$, respectively; this does not remove the hidden-constant caveat.

\subsection{End-to-end recovery from an extension-field projective slice}

The new slice does not reveal a coordinate column of the graph matrix.  Instead it gives a general oil vector, for which derivative-kernel recovery is stronger.

\begin{lemma}[Recovery and descent from the projective oil point]\label{lem:descent}
Let $L_0/K$ be any finite extension.  Suppose $\Theta_F(M,R)\ne0$, and let
\[
[t^\star]\in\PP^V(L_0)
\]
be the unique point for which $x^\star=Mt^\star$ lies in the planted oil space $\PP(\cO_0)_{L_0}$.  Put
\[
C_{x^\star}=
\begin{pmatrix}P_1x^\star&\cdots&P_mx^\star\end{pmatrix}
\in L_0^{N\times m}.
\]
Then
\[
\rank C_{x^\star}=V,
\qquad
\ker C_{x^\star}^{\mathsf T}=\cO_{L_0}.
\]
The attacker can row-reduce this kernel into the public graph chart; the resulting graph matrix has entries in $K$, is fixed by the $Q$-Frobenius, and yields the planted expanded UOV trapdoor after the public pull-back.
\end{lemma}

\begin{proof}
By \cref{lem:theta-ranks}, $M$ and $R$ have full rank.  Two $L_0$-linear subspaces of vector dimensions $V+1$ and $O$ in an $N=V+O$ dimensional space always meet nontrivially.  If their intersection had vector dimension at least two, every restricted output quadric would vanish on a positive-dimensional projective oil intersection, contradicting $\Theta_F(M,R)\ne0$.  Thus the oil intersection is a unique $L_0$-rational projective point.

Smoothness of the restricted core at $t^\star$ gives ambient derivative rank at least $V$.  Total isotropy gives
\[
\cO_{L_0}\subseteq\ker C_{x^\star}^{\mathsf T},
\]
so $\rank C_{x^\star}\le N-O=V$.  Equality follows, and the two $O$-dimensional spaces coincide.  The true public oil space is defined over $K$ and transverse to the normalized public oil-coordinate factor, so its graph matrix recovered over $L_0$ (or after scalar extension to $\widehat L$) is automatically $Q$-Frobenius fixed.  Direct verification of the graph identities and total isotropy makes the procedure Las Vegas.
\end{proof}

\begin{lemma}[Projective Frobenius filter]\label{lem:projective-frobenius}
Let $k_0=\F_{q_0}$ be a finite field and let $[X_0:\cdots:X_V]\in\PP^V(\overline{k_0})$.  Then the point is $k_0$-rational if and only if
\begin{equation}\label{eq:projective-frobenius}
X_i^{q_0}X_j-X_iX_j^{q_0}=0
\qquad(0\le i<j\le V).
\end{equation}
\end{lemma}

\begin{proof}
The equations say exactly that the Frobenius image of the nonzero coordinate vector is proportional to the vector itself.  Equivalently, the associated projective point is fixed by $\operatorname{Gal}(\overline{k_0}/k_0)$, whose fixed points on projective space are $\PP^V(k_0)$.
\end{proof}

\paragraph{Root extraction.}
Corollary~5.5 returns a geometric resolution over $\widehat L$, not necessarily over the coefficient field $L_0$.  Let $q(T)\in\widehat L[T]$ be its separable degree-$2^V$ separator polynomial.  After undoing the solver's coordinate change, compute homogeneous coordinate functions
\[
X_0(T),\ldots,X_V(T)\in A:=\widehat L[T]/(q),
\]
using the inverse of $q'(T)$ in $A$ when the Kronecker output is in the usual $w_i/q'$ form.  Put $q_0=|L_0|$.  Starting from $q$, take successive gcds with polynomial representatives of
\[
X_i(T)^{q_0}X_j(T)-X_i(T)X_j(T)^{q_0}
\quad(0\le i<j\le V).
\]
By \cref{lem:projective-frobenius}, this keeps exactly the separator roots whose \emph{projective} points are $L_0$-rational.  This test remains correct even though the separator, the coordinate change, and a chosen affine representative are defined only over $\widehat L$; requiring $T\in L_0$ or requiring every displayed affine coordinate to be Frobenius fixed would in general be wrong.

Next evaluate each of the unrecombined $m-V$ equations in $A$ and take successive gcds, factor the residual polynomial over $\widehat L$, and reconstruct the corresponding projective coordinates.  An $L_0$-rational projective point has a separator value in $\widehat L$, so it contributes a linear factor there.  Fast modular exponentiation uses $O(\log q_0)$ quotient-ring products.  Frobenius filtering, gcds, factorization, derivative-kernel linear algebra, and public verification therefore cost $\widetilde O(2^V\operatorname{poly}(V,O,\log|\widehat L|))$ bit operations, below the nearly quadratic-in-$2^V$ solver bound because the internal extension degree is polynomial in $V$.  Every survivor is checked against all original forms, so filtering errors cannot produce a false trapdoor.

\subsection{Attack theorem and costs}

\begin{theorem}[Fast almost-all-key expanded-trapdoor recovery]\label{thm:attack}
For each recommended parameter set, let $L_0/K$ have degree $s_0=4,5,6$, respectively, and consider the following public trial:
\begin{enumerate}
\item sample $M\in L_0^{(V+O)\times(V+1)}$ once; if $\rank M<V+1$, return \textnormal{\textsc{fail}}, and otherwise restrict all public extension-field quadrics to $\PP(\operatorname{im}M)$;
\item sample $R\in L_0^{V\times m}$ once; if $\rank R<V$, return \textnormal{\textsc{fail}}, and otherwise run a bounded, certificate-checked invocation of the small-field Kronecker solver on the $V$ recombined restricted quadrics; the solver may construct $\widehat L/L_0$, and failed inversions, a missing a posteriori certificate, or budget exhaustion return \textnormal{\textsc{fail}};
\item use \eqref{eq:projective-frobenius} to extract the $L_0$-rational projective roots of the resolution, then retain only those satisfying every unrecombined restricted equation;
\item at each survivor, compute the ambient derivative kernel and require dimension $O$, total isotropy, Frobenius descent to $K$, and every public graph identity;
\item reconstruct the expanded central matrices and secret linear map as in \cref{prop:verified-key}.  Exercise the expanded signing algebra on a test message and accept a signature only after verification under the original public key.
\end{enumerate}
In the ideal-field experiment, the set of keys on which no pair $(M,R)$ gives a smooth full core has probability below
\[
2^{-364.56},\quad2^{-457.92},\quad2^{-1038.06}.
\]
For every key outside that set, the public slice-and-output part of one trial fails with probability at most
\[
2^{-19.85},\quad2^{-15.73},\quad2^{-9.85}.
\]
Respectively $8,14,28$ independent public trials make these terms smaller than
\[
2^{-143},\quad2^{-207},\quad2^{-272}.
\]
On every smooth trial, the planted oil point is among the filtered roots and its derivative kernel is the planted oil space.  Repeated certificate-checked algebraic trials therefore recover the planted expanded graph trapdoor with probability one on every nonexceptional key; no seed preimage is recovered or required.  On an arbitrary key, a fixed-budget version is one-sided and may report failure.

The lower-order restriction, projective rational-root filtering, derivative-kernel, descent, and public verification stages do not change the solver exponent.  At $\epsilon=0.2$, the resulting direct bit-complexity indicators, with generic circuit-simulation indicators in parentheses, are
\[
119.59\ (126.49),\qquad172.71\ (180.14),\qquad230.17\ (238.02).
\]
Running the target-amplified numbers of trials in parallel gives circuit indicators $129.49,183.95,242.83$.
\end{theorem}

\begin{proof}
The key-level exceptional probability is \cref{thm:key-density}.  On every nonexceptional key, \cref{cor:joint-sz} bounds the chance that the unconditioned draw $(M,R)$ fails to give a smooth full core, including rank failures.  Conditional on a smooth core, \cref{lem:full-implies-prefix} supplies the hypotheses of the van der Hoeven--Lecerf solver, whose a posteriori certificate makes bounded internal failures detectable.  The unique projective oil intersection, derivative-kernel recovery, and descent are \cref{lem:descent}; \cref{lem:projective-frobenius} ensures that the planted $L_0$-rational projective point survives extraction even though the resolution is over $\widehat L$; and \cref{prop:verified-key} gives the exact expanded UOV factorization.  The planted graph is therefore produced on every successful regularization trial.  All root filtering, graph tests, signing attempts, and final signature checks use the unrecombined public equations or the original public verifier, so no false output is accepted.

The post-solve operations are quasi-linear in degree $2^V$, up to polynomial factors in the parameters and in the bit size of the solver's internal field, and are lower order than the nearly quadratic solver bound.  Substituting the regularization-field degrees into \eqref{eq:fast-cost} gives the direct indicators; \eqref{eq:circuit-cost} gives the circuit-simulation indicators.  Public and internal expected retry factors are constant in the asymptotic ledger.  A parallel batch of $t$ public trials multiplies circuit size by at most $t$, producing the amplified indicators.
\end{proof}

\begin{corollary}[Fixed-budget concrete success]\label{cor:concrete-budget}
Let $t\ge1$, and run $t$ independent public slice-and-output trials over $L_0$.  Amplify the certificate-checked internal solver so that the probability that all otherwise-good trials fail internally is at most $\beta_t$.  In the ideal-field experiment, the probability of recovering no verified expanded trapdoor is at most
\[
\eta_{\off}+\eta_{\oil}+\eta_{\dep}+\delta_V(L_0)^t+\beta_t.
\]
For the concrete SHAKE/AES-expanded distribution, the same failure probability is at most the displayed quantity plus
\[
\Adv^{\rm exp}_{\QRUOV}(D_t)+\varepsilon_{\rm RS},
\]
where $D_t$ is the fixed-budget verified-success distinguisher and has the cost of those $t$ attack trials plus key generation.
\end{corollary}

\begin{proof}
Split on the event $\Theta_F\equiv0$ and apply \cref{thm:key-density}.  Outside it, independence of the public trial coins and \cref{cor:joint-sz} bound simultaneous algebraic regularization failure by $\delta_V(L_0)^t$; add the amplified internal failure $\beta_t$.  The concrete statement is \cref{prop:transfer} applied to the verified-success event with all public coins included in $U$.
\end{proof}

\subsection{Artifact-backed numerical audit}

The accompanying standard-library script \path{qruov_parameter_audit.py} evaluates the formulas used in the main theorem at high precision and contains frozen assertions for the three regularization-field degrees, the direct and circuit indicators, and the target-amplified circuit indicators.  It regenerates the key-bad bounds, joint-discriminant degrees, one-trial probabilities, expected retry factors, target and rejection-transfer amplification counts, nominal margins, rejection-transfer terms, epsilon sensitivity thresholds, and coefficient-field quotient-vector lower bounds retained in this revision.  No empirical small-field experiment is used as evidence for the main theorem, and no such experiment is claimed as part of the accompanying artifact.

\paragraph{Concrete seed expansion.}
The rejection-term logarithms in \cref{prop:transfer} are $-121.23,-184.71,-248.28$.  With one public regularization trial, the geometric failure term is larger at all three levels.  Respectively $7,12,26$ independent public trials make the geometric term smaller than the rejection-transfer term; $8,14,28$ make it smaller than the nominal category target.  The rejection term is a coupling loss in the ideal-to-concrete reduction, not an algebraic failure probability, and is therefore displayed separately.

\section{Conditional secondary direction: pseudo-oil direct forgery}\label{sec:pseudo-oil}

\begin{center}
\fbox{\begin{minipage}{0.94\textwidth}
\textbf{Status of this section.}  This direction is not part of \cref{thm:intro,thm:attack} or the primary security ledger.  It relies on the cited average-case pseudo-oil extraction and branch heuristic, and its terminal solver may use the conditional arbitrary-core appendix.  The exact algebra below begins only after a suitable decomposition has been obtained.
\end{minipage}}
\end{center}

The proved projective-slice attack recovers expanded signing data from the public extension-field tuple.  This logically independent direction acts on the full base-field verification map
\[
P:\F_{127}^{n}\longrightarrow\F_{127}^{m},
\qquad
(n,m)=(210,54),(306,78),(411,105).
\]
It uses the enriched pseudo-oil framework of May--Ostuzzi--Ressler and Ostuzzi~\cite{JustGuess,Ostuzzi}.  Those methods guess or extract a structured low-dimensional subspace on which part of the system linearizes.  The feasibility obstruction and residual elimination below are exact conditional statements; existence, extraction cost, and branch behavior remain heuristic.


\subsection{Feasibility and the obstruction after slicing}

An enriched pseudo-oil decomposition is parameterized by a guess dimension $k$, a number $\ell_{\rm po}$ of linearized equations, positive block sizes $b_1,\ldots,b_p$, and a residual dimension $r\ge\ell_{\rm po}$.  Put
\[
B=\sum_{j=1}^p b_j,
\qquad
L_{\rm po}=B+\ell_{\rm po},
\qquad
r=m-B-k.
\]
The two feasibility inequalities in the adopted average-case model are
\begin{align}
 n-(r-1)(L_{\rm po}+1)&\ge L_{\rm po},\label{eq:po-one}\\
 n-m+1&\ge\sum_{j=1}^p b_j\left(m-\sum_{i=1}^j b_i-k\right).\label{eq:po-two}
\end{align}

\begin{theorem}[Post-slice pseudo-oil obstruction]\label{thm:po-obstruction}
For an enriched pseudo-oil decomposition with $\ell_{\rm po}\ge1$:
\begin{enumerate}
\item if $n<m$, no parameters satisfy \eqref{eq:po-two};
\item if $n=m$, feasibility forces $B=\ell_{\rm po}=1$, $p=1$, $b_1=1$, and $k=m-2$.
\end{enumerate}
Thus an overdetermined projective or affine QR-UOV slice admits no useful enriched pseudo-oil step, while a square slice admits only a degenerate $q^{m-2}$ guessing attack.
\end{theorem}

\begin{proof}
Every factor in the right side of \eqref{eq:po-two} is at least $r\ge\ell_{\rm po}\ge1$, so the right side is at least $B\ell_{\rm po}>0$.  If $n<m$, the left side is nonpositive.  If $n=m$, the left side is one, forcing $B=\ell_{\rm po}=1$ and hence $p=b_1=1$; equality then gives $m-b_1-k=1$.
\end{proof}

The extension-field key-recovery tuple is underdetermined, but not enough to help this framework: \eqref{eq:po-two} implies $Br\le n-m+1$ and therefore $k\ge2m-n-2$.  At the recommended shapes this means at least $36,52,71$ guessed extension-field variables, costing roughly $2^{755},2^{1090},2^{1489}$ before the terminal solve.  The pseudo-oil step must therefore precede slicing and act on the full base-field map.

\subsection{The working pre-slice construction}

We now fix a concrete feasible decomposition. Choose
\[
p=2,
\qquad
(b_1,b_2)=(1,1),
\qquad
\ell_{\rm po}=1,
\qquad
k=0.
\]
Then $B=2$, $L_{\rm po}=3$, and $r=m-2$.  For the three recommended shapes, the left side of \eqref{eq:po-one} is respectively $6,6,3$, while the right side is $3$; \eqref{eq:po-two} has slack $52,76,100$.  Thus the adopted heuristic predicts a valid decomposition at every level.  The subspace-extraction step (called \emph{centrifugation} in the pseudo-oil framework~\cite{JustGuess,Ostuzzi}) constructs
\[
U=U_0\oplus U_1\oplus\widetilde U,
\qquad
\dim U_0=m-2,
\quad
\dim U_1=\dim\widetilde U=1,
\]
with the following pattern for the first three equations:
\begin{enumerate}[label=(\alph*)]
\item $U_0$ is totally isotropic for equations $1,2,3$;
\item $U_0\oplus U_1$ is totally isotropic for equation $1$;
\item $U_0$ is orthogonal to $U_1\oplus\widetilde U$ for equations $1,2$;
\item $U_0\oplus U_1$ is orthogonal to $\widetilde U$ for equation $1$.
\end{enumerate}

\begin{theorem}[Exact residual reduction]\label{thm:po-reduction}
Conditioned on a successful subspace extraction with the parameters above, direct forgery reduces exactly to:
\begin{enumerate}
\item one univariate quadratic in the coordinate of $\widetilde U$;
\item one univariate quadratic in the coordinate of $U_1$;
\item one affine-linear equation eliminating one $U_0$ coordinate;
\item $m-3$ quadratic equations in $m-3$ variables.
\end{enumerate}
A solution of the final square system reconstructs a preimage of the original verification target by linear algebra.
\end{theorem}

\begin{proof}
For equation one, $U_0\oplus U_1$ is totally isotropic and orthogonal to $\widetilde U$, so after fixing the target the equation is univariate quadratic in the coordinate of $\widetilde U$.  For equation two, $U_0$ is both totally isotropic and orthogonal to $U_1\oplus\widetilde U$; after the first coordinate is fixed, the equation is univariate quadratic in the coordinate of $U_1$.  For equation three, total isotropy of $U_0$ removes every $U_0$--$U_0$ term, so after the two one-dimensional coordinates are fixed the equation is affine-linear on $U_0$.  Eliminating one coordinate leaves the final $m-3$ equations on an affine space of the same dimension.  Every step is an exact substitution, so reversing the substitutions reconstructs a preimage.
\end{proof}

Under the published extraction and branch heuristic, the two univariate stages generate one partial branch in expectation and the terminal residual is modeled as a generic square system with one nonsingular isolated solution in expectation.  Failed or singular branches can be discarded, the decomposition rerandomized, or the public salt changed.  These are the same average-case assumptions as in the underlying pseudo-oil analyses; the algebraic reduction after a valid decomposition is exact.

\subsection{Costs}

The primary pseudo-oil line uses the field-uniform arbitrary-core solver from \cref{app:fallback} on the final $(m-3)$-variable system and therefore inherits the preceding nonsingular-terminal-branch assumption.  For the selected two one-dimensional blocks, the charged subspace-extraction work is $(m-2)\Solve(3,3)+\Solve(1,1)$ plus polynomial-time linear algebra; the constant-size branch tree is also lower order than the terminal solve in the adopted model.

\begin{table}[H]
\centering
\small
\setlength{\tabcolsep}{3pt}
\caption{Pseudo-oil direct forgery.  Entries are direct indicators, with uniform indicators in parentheses.}
\label{tab:pseudo-oil}
\begin{tabular}{@{}lrrr@{}}
\toprule
Level & Final & Field-uniform Strassen & Target\\
\midrule
I   & 51  & $129.52\ (136.53)$ & 143\\
III & 75  & $180.06\ (187.55)$ & 207\\
V   & 102 & $236.18\ (244.06)$ & 272\\
\bottomrule
\end{tabular}
\end{table}

\section{Security consequences, cost conventions, and limitations}\label{sec:security}

\subsection{Cost conventions}

The \QRUOV{} specification compares attacks to category targets using a Boolean-gate model~\cite{QRUOVSpec}.  Our \emph{direct} column is the base-two exponent in the soft-$O$ bit-complexity bound obtained by specializing the cited small-field solver corollary, including the input-field bit degree $3s_0$ of the regularization coefficient field.  The corollary's internally constructed extension is already included in that bound.  Our \emph{circuit} column applies the generic statement that a $T$-step bit computation admits a Boolean circuit of size $O(T\log T)$~\cite{PippengerFischer}.  Thus it adds $\log_2 C$ to a direct exponent $C$.

This conversion is an asymptotic upper-bound device, not a calibrated gate estimate.  Both columns suppress constants, soft factors, memory traffic, and the distinction between a theorem-level machine model and an optimized implementation.  Random choices can be treated as circuit inputs, a constant number of attempts can be run in parallel, and every candidate can be checked publicly.

\subsection{Primary ledger}

\begin{table}[H]
\centering
\small
\setlength{\tabcolsep}{4pt}
\caption{Primary security ledger for the proved projective-slice attack at $\epsilon=0.2$.  ``Direct'' and ``Circuit'' are asymptotic exponent indicators as defined above.  Conditional pseudo-oil and arbitrary-core estimates are deliberately excluded.}
\label{tab:master}
\begin{tabular}{@{}lrrrrrr@{}}
\toprule
Level & $[L_0:K]$ & Target & Spec. best & Direct & Circuit & Margin\\
\midrule
I   & 4 & 143 & 150 & 119.59 & 126.49 & 16.51\\
III & 5 & 207 & 211 & 172.71 & 180.14 & 26.86\\
V   & 6 & 272 & 277 & 230.17 & 238.02 & 33.98\\
\bottomrule
\end{tabular}
\end{table}

The ideal-field exceptional-key probabilities are far below the nominal targets.  One public regularization trial already succeeds with probability greater than $0.9989$ on every nonexceptional key.  The per-trial geometric terms and exact amplification counts are given in \cref{cor:joint-sz}; concrete fixed-budget success additionally carries the expansion-distinguishing and rejection-sampling terms from \cref{prop:transfer}.  These are logically different quantities and are not folded into the one-trial complexity indicators.  A target-negligible public-failure budget adds at most $3.00,3.81,4.81$ bits by the naive parallel charge and still remains below target.

\subsection{Parameter consequences}

The proved attack is governed primarily by $V=v/3$: its square core has B\'ezout degree $2^V$.  Increasing $m$ improves the key-density exponent through $e=m-V$ but does not change that leading path count.  A replacement parameter search must therefore raise $V$ enough to clear the target after charging the regularization coefficient field and a declared solver $\epsilon$, while separately rechecking key size, signing failure, and every attack already listed by the specification.  The conditional pseudo-oil section suggests that the full base-field output dimension $m$ may create a second surface, but no parameter recommendation in this paper relies on that heuristic line.

A structural countermeasure would have to remove or materially alter the public extension-field UOV tuple used by projective slicing.  Changing the public oil chart or hiding coordinate directions does not address the attack: it samples arbitrary projective planes and recovers the oil space from an ambient derivative kernel.  Within the current design, parameter inflation is the immediate response suggested by the proved analysis.

\subsection{Asymptotic, not practical}

Even one degree-$2^V$ quotient vector over the \emph{coefficient field} $L_0$ already occupies approximately $0.0409$ EiB, $8.59\times10^5$ EiB, and $6.92\times10^{13}$ EiB at Levels~I, III, and V.  These are lower bounds: the solver's internal extension increases the coefficient size, and materializing all coordinate polynomials adds a factor on the order of $V$.  No implementation in this work approaches full-parameter time or memory.  The result is an asymptotic parameter shortfall in the specification's attack model, not a practical key-recovery demonstration.

\section{Conclusion}\label{sec:conclusion}

A random projective $V$-plane necessarily meets the hidden oil space.  When a random square core on that plane is smooth, the intersection is one point and the ambient derivative kernel at that point is the full oil space.  The joint discriminant and off-oil incidence argument make this an explicit almost-all-key theorem, with ideal-field bad-key probabilities below $2^{-364.56},2^{-457.92},2^{-1038.06}$.

The technically sharper solver interface separates the public regularization field $L_0$ from the solver's internally constructed extension $\widehat L$.  Corollary~5.5 of van der Hoeven--Lecerf then yields direct bit-complexity indicators $119.59,172.71,230.17$ and generic circuit-simulation indicators $126.49,180.14,238.02$ at $\epsilon=0.2$.  The latter remain below the nominal targets by $16.51,26.86,33.98$ bits.  One public trial succeeds with probability at least $1-2^{-9.85}$ even at Level~V; making the public algebraic failure category-negligible by $8,14,28$ parallel trials leaves margins $13.51,23.05,29.17$ bits under the same asymptotic-indicator convention.

The rational-root stage must be projective because the solver's primitive element, coordinate change, and affine normalization may live only in $\widehat L$.  Frobenius proportionality of homogeneous coordinate vectors gives an exact $L_0$-rationality test and preserves the planted oil point.  The recovered object is an expanded UOV signing trapdoor rather than a seed-encoded secret key.  The pseudo-oil and arbitrary-core directions remain conditional and are not needed for the main conclusion.  All numerical rows remain asymptotic specializations with hidden constants and soft factors, not a practical full-parameter implementation.

\begin{thebibliography}{99}

\bibitem{QRUOVSpec}
H.~Furue, Y.~Ikematsu, F.~Hoshino, T.~Takagi, H.~Kosuge, K.~Yamakoshi, R.~Akiyama, S.~Nakamura, S.~Orihara, and K.~Kinjo.
\newblock \emph{QR-UOV Specification Document (Round 2 submission)}.
\newblock NIST Additional Digital Signatures Project, 5 February 2025.

\bibitem{QRUOVRef}
QR-UOV team.
\newblock \emph{Round 2 Reference Implementation of QR-UOV}.
\newblock Public source repository, December 2024.

\bibitem{KPG1999}
A.~Kipnis, J.~Patarin, and L.~Goubin.
\newblock Unbalanced Oil and Vinegar signature schemes.
\newblock In \emph{EUROCRYPT 1999}, pp.~206--222, 1999.

\bibitem{QRUOV2021}
H.~Furue, Y.~Ikematsu, Y.~Kiyomura, and T.~Takagi.
\newblock A new variant of unbalanced oil and vinegar using quotient ring: QR-UOV.
\newblock In \emph{ASIACRYPT 2021, Part IV}, pp.~187--217, 2021.

\bibitem{NISTRound3}
National Institute of Standards and Technology.
\newblock Nine candidates advance to the third round of the Additional Digital Signatures for the PQC standardization process.
\newblock NIST CSRC, 14 May 2026.

\bibitem{NISTIR8610}
G.~Alagic, M.~Bros, P.~Ciadoux, et al.
\newblock \emph{Status Report on the Second Round of the Additional Digital Signature Schemes for the NIST Post-Quantum Cryptography Standardization Process}.
\newblock NISTIR~8610, May 2026.

\bibitem{Benoist}
O.~Benoist.
\newblock Degr\'es d'homog\'en\'eit\'e de l'ensemble des intersections compl\`etes singuli\`eres.
\newblock \emph{Annales de l'Institut Fourier}, 62(3):1189--1214, 2012.

\bibitem{vdHL}
J.~van der Hoeven and G.~Lecerf.
\newblock On the complexity exponent of polynomial system solving.
\newblock \emph{Foundations of Computational Mathematics}, 21(1):1--57, 2021.

\bibitem{SafeyElDinSchost}
M.~Safey El Din and E.~Schost.
\newblock Bit complexity for multi-homogeneous polynomial system solving---application to polynomial minimization.
\newblock \emph{Journal of Symbolic Computation}, 87:176--206, 2018.

\bibitem{Schost}
E.~Schost.
\newblock Computing parametric geometric resolutions.
\newblock \emph{Applicable Algebra in Engineering, Communication and Computing}, 13(5):349--393, 2003.

\bibitem{Vu}
T.~X.~Vu.
\newblock A symbolic homotopy algorithm for solving composable polynomial systems.
\newblock In \emph{ISSAC 2026}, ACM, 2026; arXiv:2605.22514.

\bibitem{Harvey}
D.~Harvey, J.~van der Hoeven, and G.~Lecerf.
\newblock Faster polynomial multiplication over finite fields.
\newblock \emph{Journal of the ACM}, 63(6):52:1--52:23, 2017.

\bibitem{PippengerFischer}
N.~Pippenger and M.~J.~Fischer.
\newblock Relations among complexity measures.
\newblock \emph{Journal of the ACM}, 26(2):361--381, 1979.

\bibitem{Strassen}
V.~Strassen.
\newblock Gaussian elimination is not optimal.
\newblock \emph{Numerische Mathematik}, 13:354--356, 1969.

\bibitem{Beullens}
W.~Beullens.
\newblock Improved cryptanalysis of UOV and Rainbow.
\newblock In \emph{EUROCRYPT 2021, Part I}, pp.~348--373, 2021.

\bibitem{FurueIkematsu}
H.~Furue and Y.~Ikematsu.
\newblock A new security analysis against MAYO and QR-UOV using rectangular MinRank attack.
\newblock In \emph{IWSEC 2023}, pp.~101--116, 2023.

\bibitem{Suzuki}
T.~Suzuki, H.~Furue, T.~Ito, S.~Nakamura, and S.~Uchiyama.
\newblock An extended rectangular MinRank attack against UOV and its variants.
\newblock Cryptology ePrint Archive, Report 2025/739, 2025.

\bibitem{Pebereau}
P.~P\'ebereau.
\newblock One vector to rule them all: Key recovery from one vector in UOV schemes.
\newblock In \emph{PQCrypto 2024}, pp.~92--108, 2024; ePrint 2023/1131.

\bibitem{JustGuess}
A.~May, M.~Ostuzzi, and H.~Ressler.
\newblock Just Guess: Improved (quantum) algorithm for the underdetermined MQ problem.
\newblock In \emph{EUROCRYPT 2026, Part IV}, pp.~334--362, 2026; ePrint 2025/1788.

\bibitem{Ostuzzi}
M.~Ostuzzi.
\newblock Pseudo-oil subspaces and the geometry of underdetermined MQ problems.
\newblock Cryptology ePrint Archive, Report 2026/1122, 2026.

\bibitem{Jin}
Y.~Jin, Y.~Pan, X.~He, B.~Gong, and J.~Ding.
\newblock Security analysis on UOV families with odd characteristics: Using symmetric algebra.
\newblock In \emph{PKC 2026, Part I}, pp.~428--454, 2026; ePrint 2025/1137.


\bibitem{Aulbach}
T.~Aulbach, F.~Campos, and J.~Kr\"amer.
\newblock SoK: On the physical security of UOV-based signature schemes.
\newblock In \emph{PQCrypto 2025}, pp.~199--231, 2025; ePrint 2024/1818.

\bibitem{Heintz}
J.~Heintz.
\newblock Definability and fast quantifier elimination in algebraically closed fields.
\newblock \emph{Theoretical Computer Science}, 24:239--277, 1983.

\bibitem{LachaudRolland}
G.~Lachaud and R.~Rolland.
\newblock On the number of points of algebraic sets over finite fields.
\newblock \emph{Journal of Pure and Applied Algebra}, 219(11):5117--5136, 2015.

\end{thebibliography}

\appendix

\section{Attack algorithms}\label{app:algorithms}

\paragraph{Fast projective-slice expanded-trapdoor recovery.}
\begin{enumerate}
\item Convert the compressed public key to the specification's extension-field tuple over $K=\F_{127^3}$.  Choose the least $s_0$ with $Q^{s_0}>3V(V+1)2^{V-1}$ and construct $L_0=\F_{127^{3s_0}}$.
\item Draw $M\in L_0^{(V+O)\times(V+1)}$ once.  Return \textsc{fail} if $\rank M<V+1$; otherwise form all restricted quadrics $q_i(t)=t^{\trans}M^{\trans}P_iMt$.
\item Draw $R\in L_0^{V\times m}$ once.  Return \textsc{fail} if $\rank R<V$; otherwise form $g=Rq$ and run a bounded certificate-checked invocation of the small-field Kronecker solver.  It may construct $\widehat L/L_0$; abort on a failed inversion, missing certificate, or budget exhaustion.
\item Undo the solver's coordinate change.  In $\widehat L[T]/(q)$ impose all projective Frobenius relations \eqref{eq:projective-frobenius} for $q_0=|L_0|$, then impose every unrecombined restricted equation.  Factor the residual separator polynomial over $\widehat L$ and reconstruct the surviving projective points.
\item For each survivor $x=Mt$, compute $\ker(P_1x\ \cdots\ P_mx)^{\trans}$.  Require dimension $O$, Frobenius descent to $K$, total isotropy, and every public graph identity.
\item Reconstruct $F_i'$ and $S_{G'}$ as in \cref{prop:verified-key}, apply the public pull-back to expanded base-field matrices, and exercise the expanded signing algebra.  Accept only a publicly verified signature.  On failure, continue through the candidates or begin a new public trial.
\end{enumerate}
The output is expanded signing data and a verified signature, not a seed preimage or serialized secret key.

\paragraph{Conditional pseudo-oil direct forgery.}
\begin{enumerate}
\item Expand the full base-field verification map as $m$ homogeneous quadrics in $n=v+m$ variables over $\F_{127}$.
\item Conditional on the cited extraction heuristic, obtain an enriched decomposition with $p=2$, $(b_1,b_2)=(1,1)$, $\ell_{\rm po}=1$, and $k=0$.
\item Restrict to the resulting $m$-dimensional pseudo-oil space.
\item Solve equation one in the coordinate of $\widetilde U$, then equation two in the coordinate of $U_1$, backtracking over their constant-size root tree.
\item Use equation three to eliminate one $U_0$ coordinate and solve the remaining $(m-3)$-variable square system under the solver assumption stated in \cref{sec:pseudo-oil,app:fallback}.
\item Reverse the affine substitutions and verify the reconstructed signature against the original public map.
\end{enumerate}

\section{Arbitrary-core fallback}\label{app:fallback}

\begin{center}
\fbox{\begin{minipage}{0.94\textwidth}
\textbf{Status of this appendix.}  This appendix proposes a finite-characteristic specialization and dense-quadratic cost analysis of symbolic-homotopy machinery.  It is not used by the proved main theorem or the primary security ledger.  Its correctness and all downstream pseudo-oil costs should be regarded as conditional pending independent verification of the specialization and implementation model.
\end{minipage}}
\end{center}

The main attack uses a smooth complete intersection so that the fastest Kronecker theorem applies.  The fallback instead targets only the nonsingular isolated roots of an arbitrary square core
\[
Z_{\rm ns}(f)=\{y:f(y)=0,\ \det J_f(y)\ne0\}.
\]
A full-rank planted derivative guarantees that the planted root belongs to this set, regardless of every other component and every prefix ideal.  This appendix records the finite-characteristic specialization; its success condition is incomparable with the main theorem's, and its bad-key probability is smaller.

\subsection{Finite-characteristic start fiber and separator}

\begin{lemma}[Vandermonde product start system]\label{lem:vandermonde-start}
Let $K'$ contain $2V$ distinct elements $a_1,\ldots,a_{2V}$.  Put
\[
L_a(X)=X_1+aX_2+\cdots+a^{V-1}X_V+a^V,
\qquad
g_i=L_{a_{2i-1}}L_{a_{2i}}.
\]
Then $g=0$ has exactly $C=2^V$ affine roots over $K'$, every root is simple, and a geometric resolution can be constructed in $\widetilde O(VC)$ operations.
\end{lemma}

\begin{proof}
Choose one node from each pair.  The selected $V$ linear equations have an invertible Vandermonde matrix and hence one solution; its monic polynomial has exactly the selected nodes as roots, so no unselected factor vanishes.  Different selections give different roots, and the Jacobian is an invertible Vandermonde matrix times a nonzero diagonal matrix.  Gray-code enumeration plus product-tree interpolation gives the stated cost.
\end{proof}

\begin{lemma}[Full random separator]\label{lem:separator}
If the start and target nonsingular fibers each have degree at most $C$, all well-separating conditions used by the symbolic homotopy can be expressed as nonvanishing of a random linear form on at most $C^2$ nonzero vectors.  A full random linear form over $K'$ is therefore good with probability at least $1-C^2/|K'|$.
\end{lemma}

\begin{proof}
There are at most $\binom C2$ start-point differences, $\binom C2$ target-limit differences, and $C$ nonzero leading-coefficient vectors, totaling $C^2$.  A uniform linear form annihilates each fixed nonzero vector with probability at most $1/|K'|$.
\end{proof}

\subsection{Characteristic-independent local updates}

\begin{lemma}[Newton--Hensel doubling]\label{lem:newton-doubling}
Let $A$ be a commutative ring, $F\in A[T,X_1,\ldots,X_V]^V$, and work modulo $T^{2s}$.  If
\[
F(X)\equiv0\pmod{T^s}
\quad\text{and}\quad
J_F(X)\text{ is invertible modulo }T^{2s},
\]
then
\[
X^+=X-J_F(X)^{-1}F(X)\pmod{T^{2s}}
\]
satisfies $F(X^+)\equiv0\pmod{T^{2s}}$.
\end{lemma}

\begin{proof}
The correction is divisible by $T^s$.  Taylor expansion cancels the constant and linear terms, while every remaining term contains two corrections and is divisible by $T^{2s}$.  No factorial or characteristic-dependent denominator occurs.
\end{proof}

\begin{lemma}[Primitive-element renormalization]\label{lem:primitive-renorm}
Let $B=A[T]/(T^{2s})$, let $Q\in B[U]$ be monic, and represent coordinates $X_i$ modulo $Q$.  Suppose
\[
\delta(U)=\operatorname{rem}_U(\lambda(X)-U,Q)\in T^sB[U]
\]
and write $Q'(U)\delta(U)=H(U)Q(U)+R(U)$ with $\deg_U R<\deg Q$.  Define
\[
Q^*=Q-R,
\qquad
X_i^*=\operatorname{rem}_U(X_i-X_i'\delta,Q^*).
\]
Then $U\mapsto U+\delta(U)$ gives an isomorphism $B[U]/(Q^*)\to B[U]/(Q)$, with inverse $U\mapsto U-\delta(U)$; it sends $X^*$ to $X$ and makes $\lambda(X^*)=U$.  Polynomial relations satisfied by $X$ are preserved.
\end{lemma}

\begin{proof}
Every coefficient of $\delta,H,R$ is divisible by $T^s$, so products of two such quantities vanish modulo $T^{2s}$.  First-order Taylor expansion gives
\[
P(U\pm\delta)=P(U)\pm P'(U)\delta
\]
for every polynomial $P$.  Substituting $U+\delta$ into $Q^*=Q-R$ and using $Q'\delta=HQ+R$ shows that $Q^*$ maps to zero modulo $Q$; the reverse substitution similarly maps $Q$ to zero modulo $Q^*$.  The two substitutions are inverse because $\delta(U+\delta)=\delta(U)$.  The coordinate and primitive-element identities follow from the same first-order expansion, and an isomorphism preserves all polynomial relations.
\end{proof}

\begin{theorem}[Finite-characteristic nonsingular-root solver]\label{thm:fallback-solver}
Let $f$ be any $V$ affine quadrics over a finite field of odd characteristic and put $C=2^V$.  Over an extension $K'/K$ containing $2V$ distinct elements, replace the published product start system and integer-indexed separator in the Safey El Din--Schost nonsingular-solutions algorithm by \cref{lem:vandermonde-start,lem:separator}, and use the local updates above.  The output is always a subset of $Z_{\rm ns}(f)$ and contains every nonsingular isolated root with probability at least
\[
1-\frac{C^2}{|K'|}.
\]
In particular, $|K'|\ge8C^2$ gives success probability at least $7/8$.
\end{theorem}

\begin{proof}
Safey El Din--Schost's Proposition~5 assumes that the characteristic is at least the maximum of two explicit quantities.  Their proof uses these two hypotheses separately.  The first is the hypothesis of their Lemma~7 and is used to make the integer-indexed affine factors in the start system distinct; \cref{lem:vandermonde-start} supplies a simple start fiber of the same degree without that restriction.  The second is the hypothesis of their Lemma~14 and is used only to guarantee a sufficiently large integer-indexed separator family; \cref{lem:separator} replaces it by a full random linear form and the same $C^2$ bad-vector count.

The intervening target-fiber analysis explicitly passes to generalized power series in arbitrary characteristic.  After the two replacements, the published Lemma~13 proof reuses the same global lift, target specialization, squarefree extraction, and Jacobian \textsc{Clean} steps; its degree and precision quantities depend only on the unchanged start and target degrees.  The local Newton and primitive-parameter identities are supplied by \cref{lem:newton-doubling,lem:primitive-renorm}.  Thus the proof of Proposition~5 carries over with failure probability $C^2/|K'|$.  Remark~15 and the final \textsc{Clean} call give one-sidedness: a bad separator may omit roots or report failure but cannot output a point outside $Z_{\rm ns}(f)$~\cite[Lemmas~7, 13, 14, Proposition~5, and Remark~15]{SafeyElDinSchost}.
\end{proof}

\subsection{A coordinate-slice batch}

The fallback's coverage depends on how many of the $O$ coordinate directions give a nonsingular planted slice. For $1\le j\le O$, fixing the final public coordinates to $e_j$ leaves a $V$-variable affine system with planted point $(-Ge_j,e_j)$. In central coordinates and after translation it is
\[
f_{i,j}(y)=y^{\trans}A_iy+2y^{\trans}B_ie_j.
\]
Let $J_j$ be the derivative of all $m$ equations at the planted point.

\begin{theorem}[Full-derivative coordinate batch]\label{thm:coordinate-batch}
In the ideal-field experiment, $J_1,\ldots,J_O$ are independent uniform $m\times V$ matrices. Put
\[
\tau_{m,V}=1-\prod_{r=0}^{V-1}(1-Q^{r-m}).
\]
The exact probability that every direction has deficient full derivative is $\tau_{m,V}^{O}$.  Its base-two logarithms at Levels~I, III, and V are
\[
-1132.17,\qquad -1635.35,\qquad -2935.25.
\]
\end{theorem}

\begin{proof}
The $i$th row of $J_j$ is $2(B_ie_j)^{\trans}$. For fixed $i$, these are the independent columns of the uniform matrix $B_i$; independence over $i$ completes the claim.
\end{proof}

\subsection{Dense-quadratic batching and costs}

At precision $s$, the lift works in a coefficient algebra
\[
\mathcal A_s=K'[U,T]/(q(U),T^s).
\]
Flatten the $V$ quadratic Hessians into a $V^2\times V$ matrix and the lifted coordinates into a $V\times\dim\mathcal A_s$ table.  Their product forms every Jacobian coefficient simultaneously.  Any bilinear matrix-multiplication identity over $K'$ remains valid over the associative $K'$-algebra $\mathcal A_s$.  The lift carries the inverse Jacobian from the simple start fiber and updates it by Newton multiplication, so no division by an arbitrary coefficient-algebra element is introduced.  Once the Jacobian contractions are known, quadratic values follow from
\[
f_i(X)=\frac12X^{\trans}J_i(X)+\frac12b_i^{\trans}X+c_i.
\]

\begin{theorem}[Arbitrary-core lifting bound]\label{thm:fallback-cost}
With $C=2^V$ and target precision $C_0=2^{V-1}(V+2)$, the finite-characteristic solver of \cref{thm:fallback-solver} can be implemented in
\[
\widetilde O\bigl((V^\omega+V^2)M(C)M(C_0)\bigr)
=\widetilde O(V^{\omega+1}4^V)
\]
operations in $K'$.
\end{theorem}

\begin{proof}
At precision $s$, coefficient batching computes all Jacobian contractions in $\widetilde O(V^\omega M(C)M(s))$ operations.  The quadratic-value identity above costs $O(V^2M(C)M(s))$, and the inverse-Jacobian, coordinate-correction, and primitive-parameter updates fit the same bound.  The precision doubles geometrically, so superlinearity of polynomial multiplication bounds the sum by the final stage $s=C_0$.  Start parametrization, target specialization, squarefree extraction, gcd operations, and \textsc{Clean} are polynomial in $V$ times $C$ and the dense quadratic evaluation length and are dominated by the displayed $4^V$ term.  Since $M(C)=\widetilde O(2^V)$ and $M(C_0)=\widetilde O(V2^V)$, the second form follows; see also the finite-field multiplication bounds of Harvey--van der Hoeven--Lecerf~\cite{Harvey}.
\end{proof}

Combining this solver with the coordinate batch in \cref{thm:coordinate-batch} gives the following robust line.  The Strassen column is field-uniform and is regenerated by the accompanying script directly from \cref{thm:fallback-cost}.  Replacing Strassen by ordinary cubic matrix algebra gives the final uniform column.

\begin{table}[H]
\centering
\small
\setlength{\tabcolsep}{4pt}
\caption{Arbitrary-core fallback.  The Strassen entry is the direct indicator with the uniform indicator in parentheses; the cubic entry is the uniform indicator obtained from the same formula with $\omega=3$.}
\label{tab:fallback-costs}
\begin{tabular}{@{}lrrr@{}}
\toprule
Level & Field-uniform Strassen & Cubic uniform & Target\\
\midrule
I   & $131.79\ (138.83)$ & 139.91 & 143\\
III & $182.26\ (189.77)$ & 190.96 & 207\\
V   & $236.18\ (244.06)$ & 245.34 & 272\\
\bottomrule
\end{tabular}
\end{table}

Let $G$ be the number of full-rank coordinate directions.  In the ideal-field experiment $G\sim\operatorname{Bin}(O,1-\tau_{m,V})$.  Conditional on $G>0$, one trial succeeds in choosing a regular direction, an invertible square recombination, and a good separator with probability at least
\[
\frac{G}{O}
\left(\prod_{r=0}^{V-1}(1-Q^{r-V})\right)\frac78.
\]
Using $1/G\le2/(G+1)$ and
\[
\mathbb E\!\left[\frac1{G+1}\right]
=\frac{1-\tau_{m,V}^{O+1}}{(O+1)(1-\tau_{m,V})}
\]
shows that the expected number of full solver calls, averaged over ideal-field keys conditional on $G>0$, is below $16/7+o(1)$.  The exact all-directions-deficient probability is $\tau_{m,V}^O$ by \cref{thm:coordinate-batch}; fixed-budget concrete success transfers by \cref{prop:transfer}.  On every key with $G>0$, repeated public trials are Las Vegas after substitution, graph completion, isotropy, and expanded-decomposition and public-signature verification; on an all-deficient key the procedure may fail to terminate.

Under this appendix's own asymptotic ledger, the fallback circuit indicators are $12.34,9.63,6.04$ bits above the revised main line.  Because the finite-characteristic specialization is conditional, these figures are reported only as a robustness direction and not as a second proved attack.

\section{Negative attack audit}\label{app:negative}

\paragraph{Why the old one-slice exceptional-locus program was abandoned.}
For a fixed planted slice, the exceptional locus contains the determinantal set on which the full planted derivative has rank below $V$: every recombination is then singular at the planted point.  This component has codimension $m-V+1=3,3,4$ and exact ideal-field density $\tau_{m,V}$, about $2^{-62.90},2^{-62.90},2^{-83.86}$.  Thus a category-negligible density theorem for one unconditioned slice is false.  The projective attack changes the randomization space and proves a different joint exceptional-set bound.

\paragraph{All graph columns at once.}
The full coupled graph system is heavily overdetermined, but it is not a random overdefined MQ instance.  Once one graph column is known, all others are linear.  The global path-count theorem \cref{thm:bezout} shows that every nonzero multihomogeneous B\'ezout count for a square selection is at least $2^V$, attained by one-column recovery and tree-structured completion.  Equation selection alone therefore does not beat the main square solve.

\paragraph{Using all spare slice equations directly.}
The fixed affine systems have shapes $54$-in-$52$, $78$-in-$76$, and $105$-in-$102$.  Their standard semiregular Hilbert-series screens first become nonpositive at degrees $25,36,46$, with logarithmic monomial counts $66.653,97.829,128.506$.  Even an optimistic quadratic sparse-linear-algebra charge is above the fast discriminant line before field arithmetic, matrix construction, memory, and circuit conversion.

\paragraph{Conditioning and low-degree lifts.}
Exact affine Macaulay calculations through six variables agree with matched random controls at every preterminal degree; the planted point changes only the final one-dimensional quotient.  Higher tensor and moment lifts contain exact exterior-power kernels, but their row counts substantially overpredict rank because of Macaulay syzygies.  No stable recursive radical, module, or low-rank decomposition was found in the included small-dimensional controls.  These tests rule out the implemented shortcuts, not all conceivable algorithms.

\section{A determinant-pencil characterization}\label{app:pencil}

We record a structural observation linking common quadric roots to the singular locus of the determinant pencil.  It clarifies the geometry but did not lead to a faster attack.

Let $P_1,\ldots,P_m\in\Sym_N(K)$ be homogeneous sliced matrices and define
\[
M(\lambda)=\sum_i\lambda_iP_i,
\qquad
\Delta(\lambda)=\det M(\lambda).
\]
For a common root $x$, put $K_x=\ker[P_1x\ \cdots\ P_mx]$.

\begin{proposition}[Critical-locus duality]\label{prop:pencil}
The projective space $\PP(K_x)$ is contained in the singular locus of $\Delta=0$.  Conversely, every corank-one critical pencil member has a kernel vector that is a common quadric root.
\end{proposition}

\begin{proof}
For $\lambda\in K_x$, $M(\lambda)x=0$.  At corank one, $\operatorname{adj}M(\lambda)=cxx^{\trans}$, so Jacobi's identity gives
\[
\frac{\partial\Delta}{\partial\lambda_i}
=\operatorname{tr}(\operatorname{adj}M(\lambda)P_i)
=cx^{\trans}P_ix=0.
\]
At corank at least two, the adjugate is zero.  Conversely, if a critical member has corank one with kernel $y$, the same identity gives $y^{\trans}P_iy=0$ for every $i$.
\end{proof}

Every corank-at-least-two member is automatically critical and creates a large extraneous component.  Exact finite-field tests confirm the characterization.

\section{Reproducibility and scope}\label{app:reproducibility}

This revision is accompanied by the manuscript source and the standard-library script \path{qruov_parameter_audit.py}.  The script recomputes the main theorem's parameter, probability, regularization-field, complexity-indicator, amplification, rejection-transfer, sensitivity, and memory-lower-bound figures and contains frozen assertions against accidental formula drift.  It does not implement full-parameter polynomial-system solving, key recovery, or signing.

\paragraph{Exact results used by the main theorem.}
The public extension-field reduction; the cached ideal-expansion factorization in \cref{prop:ideal-factorization}; projective oil intersection; joint-discriminant nonvanishing outside the three explicit loci; finite-field incidence and degree bounds; full-core solver regularity; projective Frobenius root filtering; derivative-kernel recovery and descent; exact expanded UOV reconstruction; oil-space uniqueness codimension; and the multihomogeneous path floor.

\paragraph{Distributional or computational-transfer results.}
Key-density bounds use the ideal-field experiment.  Concrete fixed-budget verified-success experiments transfer through \cref{prop:transfer}, which adds the expansion-distinguishing advantage and rejection-sampling term.  The latter is a reduction loss rather than an experimentally measured failure rate.

\paragraph{Conditional or asymptotic components.}
The direct, circuit, and amplified-circuit columns are specializations of asymptotic bounds with hidden constants, $\epsilon$-dependent constants, and soft factors.  Pseudo-oil existence, extraction, and branch behavior are heuristic.  The finite-characteristic arbitrary-core specialization and all cost rows depending on it are conditional and are excluded from the main theorem and primary ledger.

\paragraph{Outside scope.}
The artifact does not provide a byte-exact attack on an official compressed-key known-answer test, a full-parameter execution, a practical-memory geometric-resolution implementation, calibrated constants for fast arithmetic, a seed-preimage recovery algorithm, or a theorem for malformed public keys.

\end{document}
